% pick6-engine.tex — the pick6 fantasy-football draft engine, written up as a paper.
%
% BUILD:  latexmk -pdf docs/pick6-engine.tex
%   (or:  pdflatex docs/pick6-engine.tex  — run it twice, for the cross-references)
%   `make paper` from the repo root does the same thing.
%
% Only stock TeX Live packages. No shell-escape, no external data, no bibliography
% program: the references are a hand-written thebibliography block on purpose, so
% this file compiles anywhere pdflatex does.
%
% House style note: the repo's "everything lowercase" rule is about the TUI and the
% command output. This is a paper; it uses ordinary sentence case.

\documentclass[10pt,a4paper]{article}

\usepackage[T1]{fontenc}
\usepackage[utf8]{inputenc}
\usepackage{lmodern}
\usepackage{microtype}
\usepackage{amsmath,amssymb,amsthm}
\usepackage{booktabs}
\usepackage{array}
\usepackage{longtable}
\usepackage[top=0.85in,bottom=0.9in,left=0.95in,right=0.95in]{geometry}
\usepackage{xcolor}
\usepackage[hidelinks]{hyperref}

\theoremstyle{plain}
\newtheorem{theorem}{Theorem}[section]
\newtheorem{proposition}[theorem]{Proposition}
\newtheorem{lemma}[theorem]{Lemma}
\newtheorem{corollary}[theorem]{Corollary}
\theoremstyle{definition}
\newtheorem{definition}[theorem]{Definition}
\newtheorem{example}[theorem]{Example}
\theoremstyle{remark}
\newtheorem{remark}[theorem]{Remark}

\newcommand{\code}[1]{\texttt{\small #1}}
\newcommand{\pos}{\mathsf{pos}}
\newcommand{\adp}{\mathrm{adp}}
\newcommand{\lo}{\mathrm{lo}}

\setlength{\parskip}{0.2em}
\setlength{\emergencystretch}{3em}

% Tighten the vertical rhythm: this paper is dense with displays, theorems and
% tables, and the class defaults leave it half again as long as it needs to be.
\setlength{\abovedisplayskip}{5pt plus 2pt minus 2pt}
\setlength{\belowdisplayskip}{5pt plus 2pt minus 2pt}
\setlength{\abovedisplayshortskip}{2pt plus 2pt}
\setlength{\belowdisplayshortskip}{3pt plus 2pt minus 2pt}
\makeatletter
\renewcommand\section{\@startsection{section}{1}{\z@}%
  {-2.4ex \@plus -0.8ex \@minus -.2ex}{1.3ex \@plus.2ex}{\normalfont\Large\bfseries}}
\renewcommand\subsection{\@startsection{subsection}{2}{\z@}%
  {-2.0ex\@plus -0.8ex \@minus -.2ex}{0.9ex \@plus .2ex}{\normalfont\large\bfseries}}
\renewcommand\subsubsection{\@startsection{subsubsection}{3}{\z@}%
  {-1.8ex\@plus -0.6ex \@minus -.2ex}{0.7ex \@plus .2ex}{\normalfont\normalsize\bfseries}}
\makeatother
\usepackage{enumitem}
\setlist{topsep=2pt, itemsep=1pt, parsep=1pt}
\usepackage{caption}
\captionsetup{font=small, skip=4pt}

\title{\textbf{What does waiting cost?}\\[0.4em]
\large The pick6 draft engine: survival, urgency, and one honest backtest}
\author{Written up for the repository \code{github.com/trisslazaj/pick6}}
\date{July 2026}

\begin{document}
\maketitle

\begin{abstract}
\noindent
A fantasy football snake draft asks one question over and over: the player in front of
you is good, but so is the next pick you will get --- what does it cost you to wait?
This paper is the complete write-up of \emph{pick6}, a terminal draft assistant that
answers that question numerically and shows its working. The engine models each
player's availability as a logistic survival curve through the market's average draft
position, conditions it on the fact that he is demonstrably still on the board right
now, and corrects the resulting probabilities with a single scalar so that the number
of players the model expects to disappear equals the number of picks that will
actually happen. From those probabilities it computes the expected value of the best
player who will still be there when your turn comes, and calls the difference
\emph{urgency}.

Every piece of that pipeline is derived from a definition rather than asserted, and
every piece is graded. We backtest against one real completed 12-team half-PPR draft
using era-correct 2024 average draft position, producing $15{,}923$ labelled
predictions over $168$ vantages. The shipped model scores a Brier of $0.0660$ and a
log-loss of $0.2222$ against $0.1022$ / $0.3580$ for a constant predictor and
$0.1111$ / $0.5159$ for the unconditional logistic. Decomposing the Brier score shows
where the gain came from: the reliability term falls from $0.0231$ to $0.0047$ while
resolution barely moves, which is to say the work bought \emph{calibration}, not
\emph{discrimination}.

Three ideas were built, measured, and did not survive, and they get a section of their
own, because the reasons they failed are more interesting than the reasons the survivors
worked. The evidence for the headline change rests on a single draft, which we say
plainly, repeatedly, and in the tool's own output.
\end{abstract}

\setcounter{tocdepth}{1}
\tableofcontents
\newpage

%=======================================================================
\section{Introduction: the decision problem}
%=======================================================================

\subsection{The situation}

You are sitting at pick 27 of a twelve-team snake draft. Your next pick after this one
is pick 46. Eighteen players will come off the board in between --- picks 28 through 45,
made by nine other managers picking twice each --- and you have about ninety seconds.

In front of you are perhaps a hundred and fifty players with something to recommend
them. The best receiver available is a little better than the best running back
available. But there are sixty more receivers behind him and only a handful of running
backs, and three of the last six picks were running backs, and the tier of backs you
have been eyeing has two players left in it.

The question is not ``who is the best player available''. Your eyes can answer that.
The question is:

\begin{quote}
\emph{Which of these players will still be here in eighteen picks, and what is the
best one I could plausibly be choosing from then?}
\end{quote}

Because if the answer is ``the receiver will still be here and the running back will
not'', you take the running back, even though the receiver grades out higher. The cost
of waiting is not the same for the two positions, and that asymmetry --- not raw
player quality --- is the thing a draft is actually won and lost on.

\subsection{What the engine computes, in one paragraph}

Everything in pick6 is an answer to the question above. For each available player it
estimates $\tilde p_j$, the probability he is still on the board at your next pick.
For each position it computes $\mathbb{E}[\text{best available at my next pick}]$ from
those probabilities. It subtracts that from the value of the best player at the
position \emph{now}, multiplies by how badly your roster needs that position, and calls
the result \emph{urgency}. The board sorts by urgency. The top group is the pick.
Everything else in the tool --- the tier-cliff alarms, the run banner, the two-pick
plan, the value-over-replacement tie-break --- is either an input to that number or a
different presentation of it.

\subsection{What this document is for}

The engine was specified before it was built, and the specification asserted rather
more than it derived. This document is the derivation. It has three jobs:

\begin{enumerate}
\item \textbf{Derive, don't assert.} Where a formula in the code drops out of a
definition, the two lines that get you there are printed here. Where it does not ---
where the project made a judgement call, picked a constant by feel, or accepted a
simplification --- that is said out loud in the same voice.

\item \textbf{Introduce every symbol before using it.} A logistic distribution, an
order statistic, a hazard rate, empirical Bayes, a proper scoring rule: each of these
is defined at the point where the engine first needs it, in terms of the football, not
in terms of the next abstraction up.

\item \textbf{Map the maths to the code.} Section~\ref{sec:codemap} is a table from
every numbered equation in this paper to the Go function that implements it. If you
only read one page, read that one with the repository open beside it.
\end{enumerate}

The reader we have in mind is a strong programmer and a serious fantasy player who
likes mathematics but has not formally studied probability theory. Comfort with
algebra and with reading code is assumed. Measure theory is not; neither is any prior
exposure to survival analysis, Bayesian statistics, or forecast verification.

\subsection{The one caveat that governs everything}

The headline result --- that blending this specific league's own draft history into
the price curve improves calibration --- rests on \textbf{one draft}. There is exactly
one season for which we can obtain both a completed draft from this league and the
average draft position that the market was showing at the time. Twelve seats scored
over the same $180$ picks multiplies the \emph{count} of predictions by twelve; it does
not multiply the \emph{evidence}, because all twelve seats are scoring the same real
picks. Section~\ref{sec:threats} says this at length. The tool says it too, on every
run.

\subsection{Data provenance}
\label{sec:provenance}

Nothing in pick6 is a projection of our own. Every number describing a player is
imported:

\begin{itemize}
\item \textbf{Average draft position, and its spread.} Fantasy Football Calculator's
free API, format \code{half-ppr}, twelve teams. Each row carries \code{adp},
\code{stdev}, \code{times\_drafted}, \code{high}, \code{low} and \code{bye}. The 2026
snapshot used throughout is $203$ players over $1{,}187$ drafts (window
2026-07-24\,--\,2026-07-29); the 2024 snapshot used for the backtest is $178$ players
over $906$ drafts (window 2024-08-31\,--\,2024-09-01). Attribution is requested by the
source and given here and in the tool.

\item \textbf{Player values.} FantasyCalc's \code{redraftValue}, plus its
cross-source identifier crosswalk. Covers QB/RB/WR/TE only.

\item \textbf{Player identity, injury status, and league history.} Sleeper's
unauthenticated API: the full player dump, and this user's own completed drafts.

\item \textbf{Tiers.} A hand-transcribed rankings CSV where one covers a player;
derived from value gaps otherwise.
\end{itemize}

FFC's API is free for personal and commercial use and asks only for attribution, which is
given here, in the tool's footer, and in the repository README. The project takes no credit
for any player valuation in it.

Transformations of imported values are fair game. Inventing a player value is not, and
the one place where that temptation is strongest --- kickers and defenses, which no
source values at all --- is handled in \S\ref{sec:kdef} by anchoring onto the imported
curve rather than by making a number up.

%=======================================================================
\section{Notation and the draft}
\label{sec:notation}
%=======================================================================

\subsection{Players and values}

Let $\mathcal{P}$ be the set of players on the board. Each $j \in \mathcal{P}$ carries:

\begin{center}
\begin{tabular}{ll}
\toprule
$v_j \in \mathbb{R}_{\ge 0}$ & \emph{value}: the single number everything compares. Imported. \\
$a_j \in \mathbb{R}_{> 0}$   & \emph{average draft position} (ADP): the market's average price, in picks. \\
$\varsigma_j > 0$            & the observed standard deviation of his draft position. \\
$\sigma_j > 0$               & the \emph{scale} of his survival curve, derived from $\varsigma_j$ in \S\ref{sec:sigma}. \\
$\pos(j)$                    & his position: one of QB, RB, WR, TE, K, DEF. \\
$t_j \in \mathbb{Z}_{\ge 0}$ & his \emph{tier} within his position; $0$ means ``untiered''. \\
\bottomrule
\end{tabular}
\end{center}

\noindent
Values are only ever compared by \emph{difference}, so their absolute scale is
irrelevant --- but it must be consistent. On the shipped 2026 board the imported value
curve runs from $10{,}457$ at the top to single digits at the bottom, and it is convex:
value falls fast at the top of the board and slowly at the bottom, which is what makes
a gap between two adjacent players mean something. A linear rank would make every gap
identical and destroy the entire urgency signal.

\subsection{The snake}

A draft has $T$ teams (\emph{seats}, or \emph{slots}) and $R$ rounds, so $RT$ overall
picks numbered $1, \dots, RT$. Everything is 1-indexed, including the Go code, which
carries one $-1$ at each array access rather than an off-by-one in the arithmetic.

\begin{definition}[Round and index]
For overall pick $p \in \{1, \dots, RT\}$,
\begin{align}
r(p) &= \left\lfloor \frac{p-1}{T} \right\rfloor + 1 && \text{the round containing $p$,} \label{eq:round}\\
i(p) &= \big((p-1) \bmod T\big) + 1 && \text{the position of $p$ within its round.} \label{eq:index}
\end{align}
\end{definition}

The draft order is a permutation $\omega : \{1,\dots,T\} \to \{1,\dots,T\}$, read from
Sleeper's \code{draft\_order} field; $\omega(i)$ is the seat that picks $i$-th in an odd
round. In a \emph{snake} draft the order reverses every other round:

\begin{definition}[Seat map]
\begin{equation}
\mathrm{slot}(p) =
\begin{cases}
\omega\big(i(p)\big)      & r(p) \text{ odd},\\[0.2em]
\omega\big(T + 1 - i(p)\big) & r(p) \text{ even}.
\end{cases}
\label{eq:slotat}
\end{equation}
\end{definition}

Two facts are worth proving, not because they are hard but because getting either one
wrong corrupts every probability downstream and does so silently. A board that
mis-attributes picks still draws.

\begin{proposition}[Every pick belongs to exactly one seat]
\label{prop:bijection}
The map $p \mapsto (r(p), i(p))$ is a bijection from $\{1,\dots,RT\}$ onto
$\{1,\dots,R\} \times \{1,\dots,T\}$, and for each fixed round the seat map
$i \mapsto \mathrm{slot}$ of \eqref{eq:slotat} is a bijection onto $\{1,\dots,T\}$.
Consequently $\mathrm{slot}(\cdot)$ assigns exactly one seat to every pick.
\end{proposition}

\begin{proof}
Write $p - 1 = qT + m$ with $0 \le m < T$; the division algorithm says this
decomposition exists and is unique, and $r(p) = q+1$, $i(p) = m+1$. Ranging $p$ over
$1,\dots,RT$ ranges $q$ over $0,\dots,R-1$ and $m$ over $0,\dots,T-1$ exactly once
each, which is the first claim.

For the second: in an odd round the seat map is $\omega$ itself, a permutation. In an
even round it is $\omega \circ \rho$ where $\rho(i) = T+1-i$. The map $\rho$ is an
involution on $\{1,\dots,T\}$ (it is its own inverse, since $T+1-(T+1-i) = i$), hence a
bijection, and a composition of bijections is a bijection.
\end{proof}

\begin{corollary}[One pick per seat per round]
\label{cor:oneper}
In each round, each of the $T$ seats picks exactly once.
\end{corollary}

\begin{proof}
Immediate from the second half of Proposition~\ref{prop:bijection}: within a fixed
round, the $T$ values of $i$ map onto the $T$ seats bijectively.
\end{proof}

\subsection{My pick sequence}

Let $s \in \{1, \dots, T\}$ be \emph{my seat position} --- the index at which my seat
appears in $\omega$, so $\omega(s) = \text{mySlot}$.

\begin{proposition}[My pick in round $r$]
\label{prop:mypick}
\begin{equation}
\mathrm{myPick}(r) =
\begin{cases}
(r-1)T + s          & r \text{ odd},\\
(r-1)T + (T - s + 1) & r \text{ even}.
\end{cases}
\label{eq:mypick}
\end{equation}
\end{proposition}

\begin{proof}
In an odd round I pick when $\omega(i) = \omega(s)$, i.e.\ when $i = s$; substituting
into $p = (r-1)T + i$ gives the first line. In an even round I pick when
$\omega(T+1-i) = \omega(s)$, i.e.\ $T+1-i = s$, i.e.\ $i = T-s+1$, which gives the
second.
\end{proof}

\begin{definition}[The horizon]
Let $p_{\text{now}}$ be the current overall pick --- the next one to be made. Then
\begin{align}
q_1 &= \min\{\, \mathrm{myPick}(r) : \mathrm{myPick}(r) \ge p_{\text{now}} \,\} && \text{my next pick,} \label{eq:nextpick}\\
q_2 &= \min\{\, \mathrm{myPick}(r) : \mathrm{myPick}(r) > q_1 \,\} && \text{the one after that,} \label{eq:following}\\
N   &= q_1 - p_{\text{now}} && \text{picks until mine.} \label{eq:picksuntil}
\end{align}
$q_2 = 0$ by convention when the draft ends first. The \emph{intervening set} is
$\{p_{\text{now}}, \dots, q_1 - 1\}$, the $N$ picks that can remove a player before I
act.
\end{definition}

If it is already my turn then $q_1 = p_{\text{now}}$, $N = 0$, and the intervening set
is empty. This is not a degenerate case to be special-cased --- it is the correct
answer, and the whole engine collapses gracefully onto it: nobody can be taken in zero
picks, so every survival probability is exactly $1$, so urgency is exactly $0$
everywhere, and the board falls back to a value ordering. That is what you want on the
frame where you are actually choosing.

\subsection{The turn}

A small consequence worth deriving, because it is the geometry that makes drafting from
the edge of the snake feel different from drafting from the middle.

\begin{corollary}[The gap at the turn]
\label{cor:turn}
For odd $r$, $\mathrm{myPick}(r) + \mathrm{myPick}(r+1) = 2rT + 1$, so the gap between
my two picks around the turn is
\[
\mathrm{myPick}(r+1) - \mathrm{myPick}(r) = 2(T - s) + 1,
\]
while the gap from an even round into the next odd round is $2s - 1$.
\end{corollary}

\begin{proof}
For $r$ odd, $\mathrm{myPick}(r) = (r-1)T + s$ and $\mathrm{myPick}(r+1) = rT + (T-s+1)$
by \eqref{eq:mypick}. Adding gives $(2r-1)T + T + 1 = 2rT+1$; subtracting gives
$rT + T - s + 1 - (r-1)T - s = 2T - 2s + 1$. The even-to-odd case is the same
computation with the two lines of \eqref{eq:mypick} swapped: $rT + s$ minus
$(r-1)T + T - s + 1$ is $2s-1$.
\end{proof}

Concretely, at $T = 12$ from seat $s = 3$: the round-1 pick is $3$, the round-2 pick is
$22$ and the round-3 pick is $27$, so the gaps are $19$ and then $5$. The two gaps sum to
$2T = 24$ every time.

Do not confuse a gap with $N$. The gap is $\mathrm{myPick}(r+1) - \mathrm{myPick}(r)$ and
counts my own pick at one end; $N$ of \eqref{eq:picksuntil} counts only the picks that can
still remove somebody. Standing immediately after picking at $3$ we have
$p_{\text{now}} = 4$ and $q_1 = 22$, so $N = 18$ (picks $4$ through $21$); after picking at
$22$, $N = 4$. Gaps $19$ and $5$; horizons $18$ and $4$. \code{PicksUntilMine} prints the
second pair, and the exactly-$N$ tilt of \S\ref{sec:tilt} is solved against it.

Seat 3 lives with long waits followed by
back-to-back picks; seat 6 waits about twelve picks every time. The engine does not
need to know this --- it falls out of \eqref{eq:mypick} --- but it explains why the
two-pick lookahead of \S\ref{sec:plan} matters far more at the edges of the snake than
in the middle.

\subsection{A note on letters}

Two dozen quantities and twenty-six letters; a few get reused, so here is the whole list
of collisions up front rather than as a surprise later.

\begin{itemize}
\item $p$ is a \emph{pick number} in \S\ref{sec:notation}--\S\ref{sec:survival} and a
\emph{probability} from \S\ref{sec:tilt} onward ($p_j$, $\tilde p_j$, $\hat p$). Pick
numbers always come with a subscript that is a word ($p_{\text{now}}$) or no subscript at
all; probabilities always carry a player index or a hat.
\item $R$ is the number of \emph{rounds}, always. My remaining picks are $L$
(\S\ref{sec:urgency}) and the rounds still to play are $R_{\text{rem}}$.
\item $s$ is my \emph{seat}, always. The logistic's scale parameter is $\sigma$ (not $s$,
as most references write it), dedicated lineup slots are $\ell_P$, and the scoring rules
of \S\ref{sec:calibration} are $\mathcal{S}$.
\item Lower-case $n$ is always a \emph{count of something}: $n_j$ drafts behind a
player's spread, $n_0$ a pseudo-count, $n_k$ drafts reaching depth $k$, $n_b$ predictions
in a reliability bin, bare $n$ the total prediction count in \S\ref{sec:calibration}.
Capital $N$ is picks-until-mine and nothing else.
\item $D_j$ is the random pick at which player $j$ comes off the board;
$\mathrm{dem}(P)$, spelled out, is a position's draft demand.
\end{itemize}

\subsection{State}

The engine is a set of pure functions over one struct. It holds the players, a
$\mathrm{taken}$ set, every seat's roster so far, my seat, $T$, $R$, $p_{\text{now}}$,
$\omega$, and the league's starting-lineup shape. It recomputes everything from scratch
on every pick event. This is a few thousand floating-point operations and a handful of
bisections; there is no caching anywhere, deliberately, because a cache is a place for
two frames of the board to disagree.
%=======================================================================
\section{The survival model}
\label{sec:survival}
%=======================================================================

\subsection{What we are trying to estimate}

Fix a player $j$ and a pick number $p$. Define the event
\[
A_j(p) = \{\text{player $j$ is still undrafted immediately before pick $p$}\}.
\]
We want $\Pr[A_j(q_1) \mid A_j(p_{\text{now}})]$: given that he is sitting there right
now, what is the chance he is still sitting there when my turn comes?

Everything in this section is about turning two numbers the market gives us --- an
average price $a_j$ and a spread $\varsigma_j$ --- into that conditional probability,
without pretending to know more than those two numbers support.

\subsection{The price curve: why a logistic}

\begin{definition}[Logistic distribution]
A random variable $X$ has a \emph{logistic distribution} with location $\mu$ and scale
$\sigma > 0$ when its cumulative distribution function is
\begin{equation}
F(x) = \Pr[X \le x] = \frac{1}{1 + e^{-(x-\mu)/\sigma}}.
\label{eq:logisticcdf}
\end{equation}
It is a symmetric, bell-shaped distribution --- it looks like a normal distribution to
the eye --- with heavier tails and, crucially, a CDF in closed form.
\end{definition}

Model the pick at which player $j$ comes off the board as a logistic random variable
$D_j$ with location $a_j$ (his ADP --- the market's average price is the centre of the
distribution) and scale $\sigma_j$. Then his \emph{survival function}, the probability
he is still available at pick $p$, is
\begin{equation}
S_j(p) \;=\; \Pr[D_j \ge p] \;=\; 1 - F(p) \;=\; \frac{1}{1 + e^{(p - a_j)/\sigma_j}}.
\label{eq:S}
\end{equation}

Read it at the table. At $p = a_j$ the exponent is zero and $S_j = 1/2$: a player whose
ADP is exactly your next pick is a coin flip, which is what an average price
\emph{means}. Well before his ADP the exponent is very negative and $S_j \approx 1$.
Well after, $S_j \approx 0$.

Three reasons for the logistic over the obvious alternative, a normal (Gaussian):

\begin{enumerate}
\item \textbf{Closed form.} The normal CDF has no elementary expression; you need
\code{erf} and a numerical approximation. \eqref{eq:S} is one \code{exp} and one
division, and --- see \S\ref{sec:logspace} --- its logarithm is a function we can
evaluate stably at any magnitude. The engine calls this a few thousand times per frame
and needs it to be boring.

\item \textbf{Slightly fatter tails.} Draft position is not Gaussian: a player
occasionally goes twenty picks early because one manager loves him. Excess kurtosis
$1.2$ against the normal's $0$ is not a heavy-tailed model by any serious standard, but
it is the right direction, and it is free.

\item \textbf{The scale parameter has a clean physical reading.} This is the real
reason, and \S\ref{sec:hazard} makes it precise: deep in the tail, $\sigma_j$ becomes
a \emph{per-pick hazard} --- the chance that each passing pick is the one that takes
him. A parameter that means something at the table is worth more than a marginally
better distributional fit.
\end{enumerate}

We are not claiming draft position is really logistic. We are claiming that a
two-parameter symmetric curve through the market's mean and spread is the most a
mean and a spread can honestly support, and that of the available two-parameter
curves this is the most convenient.

\subsection{From the market's spread to the curve's scale}
\label{sec:sigma}

The source gives us $\varsigma_j$, the \emph{standard deviation} of the pick at which
$j$ was taken across its sample of real drafts. The curve wants $\sigma_j$, the scale
parameter of \eqref{eq:S}. These are not the same number, and the conversion is exact.

\begin{proposition}[Variance of a logistic]
\label{prop:logisticvar}
If $X$ is logistic with scale $\sigma$, then $\operatorname{Var}(X) = \sigma^2\pi^2/3$;
hence $\mathrm{sd}(X) = \sigma\pi/\sqrt{3}$ and
\begin{equation}
\sigma \;=\; \mathrm{sd}(X)\cdot\frac{\sqrt 3}{\pi} \;=\; \frac{\mathrm{sd}(X)}{1.8138\ldots}.
\label{eq:sigmafromstdev}
\end{equation}
\end{proposition}

\begin{proof}
Scale first: if $X$ is logistic with location $\mu$ and scale $\sigma$ then
$(X-\mu)/\sigma$ is standard logistic ($\mu = 0$, $\sigma=1$), directly from
\eqref{eq:logisticcdf}. Variance scales by $\sigma^2$ and is unaffected by location, so it
suffices to show the standard logistic has variance $\pi^2/3$.

The standard logistic's moment generating function is a classical result,
\begin{equation}
M(t) = \mathbb{E}\big[e^{tX}\big] = \Gamma(1-t)\,\Gamma(1+t) = \frac{\pi t}{\sin \pi t},
\qquad |t| < 1,
\label{eq:mgf}
\end{equation}
the second equality being Euler's reflection formula. Take logarithms to get the
cumulant generating function $K(t) = \log M(t)$, whose Taylor coefficients are the
cumulants; in particular the variance is $2\times$ the coefficient of $t^2$, since
$K(t) = \mathbb{E}[X]\,t + \tfrac{1}{2}\operatorname{Var}(X)\,t^2 + O(t^3)$.

Now do the algebra. Using $\sin u = u\left(1 - u^2/6 + O(u^4)\right)$ with $u = \pi t$,
\[
K(t) = \log(\pi t) - \log \sin(\pi t)
     = \log(\pi t) - \log\!\Big(\pi t\big(1 - \tfrac{(\pi t)^2}{6} + O(t^4)\big)\Big)
     = -\log\!\Big(1 - \tfrac{(\pi t)^2}{6} + O(t^4)\Big).
\]
Since $\log(1-x) = -x + O(x^2)$,
\[
K(t) = \frac{\pi^2 t^2}{6} + O(t^4) = \frac{1}{2}\left(\frac{\pi^2}{3}\right)t^2 + O(t^4),
\]
so $\operatorname{Var}(X) = \pi^2/3$, and $\mathbb{E}[X] = 0$ falls out of the missing
$t^1$ term (as it must, by symmetry).
\end{proof}

So the engine sets
\begin{equation}
\sigma_j \;=\; \operatorname{clamp}\!\left(\frac{\varsigma_j}{1.8138},\;
\sigma_{\min},\; \sigma_{\max}\right),
\qquad \sigma_{\min} = 0.5,\ \sigma_{\max} = 25,
\label{eq:sigmaclamp}
\end{equation}
falling back to $\sigma_{\text{def}} = 6.0$ when the source reports no spread at all.

\subsubsection*{Why a per-player scale, and not a constant}

The original specification used a single hardcoded $\sigma = 6.0$ for everybody. That
number is not wrong --- it is well calibrated for the median player. Measured across the
2026 FFC half-PPR pool (the $203$-player snapshot of \S\ref{sec:provenance}; every 2026
figure in this section is that pool, and not the PPR or standard feeds, whose spreads run
wider), $\varsigma$ runs from $0.4$ to $39.2$ with a median of $10.3$, and
$10.3/1.8138 = 5.68$. So $6.0$ is within about six percent of the median player's own
scale. But the median is not where a draft assistant earns anything:

\begin{itemize}
\item The tightest player on the board is Puka Nacua at $\varsigma = 0.4$, giving
$\sigma = 0.22$ --- clamped up to $\sigma_{\min} = 0.5$. Even after the clamp he is
\textbf{twelve times} more predictable than $\sigma = 6$ assumes. He is simply gone. Stop
hoping.
\item The widest is Evan McPherson, a round-14 kicker, at $\varsigma = 39.2$, giving
$\sigma = 21.6$. He is \textbf{three and a half times} less predictable than $\sigma = 6$
assumes. He might last three more rounds. Stop panicking.
\end{itemize}

Note what the second bullet does \emph{not} say: $\sigma_{\max} = 25$ would need
$\varsigma \ge 45.3$, and nobody on this board is near that. So of the two clamps in
\eqref{eq:sigmaclamp}, the ceiling binds on \emph{zero} of the $203$ players and the floor
on three. The 2024 era board of \S\ref{sec:backtest} reads the same $3$ low, $0$ high,
which \code{calibrate -tune} prints on every run.

Both tails are exactly where the tool is supposed to be useful, so we use the real
number where we have it. The constant survives only as the fallback for user-supplied
CSVs and any source without a spread column.

\subsection{Conditioning on the present}
\label{sec:conditioning}

Here is the single most important correction in the whole model, and it is one line of
elementary probability.

\begin{proposition}[Conditional survival]
\label{prop:conditional}
For $q \ge p$,
\begin{equation}
\Pr\big[A_j(q) \mid A_j(p)\big] \;=\; \frac{S_j(q)}{S_j(p)}.
\label{eq:psurvive}
\end{equation}
\end{proposition}

\begin{proof}
Availability is monotone: once a player is drafted he stays drafted, so for $q \ge p$
the event $A_j(q)$ \emph{implies} $A_j(p)$, i.e.\ $A_j(q) \subseteq A_j(p)$. Therefore
$A_j(q) \cap A_j(p) = A_j(q)$, and the definition of conditional probability gives
\[
\Pr[A_j(q) \mid A_j(p)] = \frac{\Pr[A_j(q) \cap A_j(p)]}{\Pr[A_j(p)]}
 = \frac{\Pr[A_j(q)]}{\Pr[A_j(p)]} = \frac{S_j(q)}{S_j(p)}. \qedhere
\]
\end{proof}

That is the whole derivation --- it is the definition of conditional probability
applied to a nested pair of events --- but its effect on the board is enormous.

\begin{example}[The faller]
A player with $a_j = 18$ and $\sigma_j = 3$ is still on the board at pick $21$, and my
next pick is $22$: one pick to go.

The \emph{unconditional} curve says
$S_j(22) = 1/\!\left(1+e^{(22-18)/3}\right) = 1/(1+3.794) = 0.209$. The board would show
$21\%$ and imply he is almost certainly gone --- across a single pick.

The \emph{conditional} ratio divides by what we already know:
$S_j(21) = 1/(1+e^{1}) = 0.269$, so
\[
\Pr[A_j(22)\mid A_j(21)] = \frac{0.2086}{0.2689} = 0.776 .
\]
Seventy-eight percent, which is right: the worst that can happen in one pick is that one
team takes him.
\end{example}

The unconditional curve is not answering our question at all. It answers ``how likely
was this player to last to pick 22, asked before the draft started'', which is a
question about the past. We already know how the past went: he is sitting there. The
ratio throws away exactly the information we have already observed and keeps exactly
the uncertainty that remains.

Conditioning barely moves anyone whose ADP is still ahead of the current pick, because
then $S_j(p_{\text{now}}) \approx 1$ and dividing by it does nothing. It only repairs
the players the market has already walked past --- and those are precisely the players
a draft assistant exists to point at.

Section~\ref{sec:results} scores this directly: the unconditional logistic
(``baseline: unconditional'') is the \emph{worst} model in the whole comparison, at
Brier $0.1111$ against $0.1022$ for a constant predictor. It is worse than not
modelling at all.

\subsection{The deep tail is a constant per-pick hazard}
\label{sec:hazard}

\begin{proposition}[Tail behaviour]
\label{prop:tail}
Fix $q > p$. As $p - a_j \to \infty$ (the player has fallen far past his price),
\begin{equation}
\frac{S_j(q)}{S_j(p)} \;\longrightarrow\; \exp\!\left(-\frac{q-p}{\sigma_j}\right).
\label{eq:tailhazard}
\end{equation}
\end{proposition}

\begin{proof}
Write $x = (p-a_j)/\sigma_j$ and $y = (q-a_j)/\sigma_j$, so $y - x = (q-p)/\sigma_j > 0$.
Then
\[
\frac{S_j(q)}{S_j(p)} = \frac{1+e^{x}}{1+e^{y}}
 = e^{x-y}\cdot\frac{e^{-x}+1}{e^{-y}+1}.
\]
As $x \to \infty$ (and hence $y \to \infty$), both $e^{-x}$ and $e^{-y}$ vanish and the
second factor tends to $1$, leaving $e^{x-y} = e^{-(q-p)/\sigma_j}$.
\end{proof}

\begin{definition}[Hazard rate]
The \emph{hazard} of an event over a small window is the probability it happens in that
window given that it has not happened yet. Here the window is one pick.
\end{definition}

Equation \eqref{eq:tailhazard} says that in the tail, surviving $k$ picks costs a factor
$e^{-1/\sigma_j}$ each time: the picks are independent trials with a constant per-pick
survival probability. The per-pick \emph{hazard} --- the chance that this pick is the
one that takes him --- is therefore
\begin{equation}
h_j \;=\; 1 - e^{-1/\sigma_j}.
\label{eq:perpick}
\end{equation}

At the table: a player with $\sigma_j = 6$ has $h_j = 15.4\%$, so roughly one pick in
six-and-a-half takes him. A player with $\sigma_j = 0.5$ has $h_j = 86.5\%$; he is not
surviving the next pick, never mind the next nineteen. That is what the scale parameter
\emph{is}, once the market has passed a player: how many picks he can absorb before
somebody bites.

\subsection{The numerics, and a genuinely instructive floating-point bug}
\label{sec:logspace}

Equation \eqref{eq:psurvive} is a ratio of two numbers that both go to zero fast. Getting
it right in floating point is not automatic, and getting it wrong produces a specific,
seductive failure that this project shipped and then measured.

\begin{definition}[Softplus]
$\operatorname{softplus}(x) = \log(1 + e^x)$. It is smooth, increasing,
$\approx 0$ for $x \ll 0$, and $\approx x$ for $x \gg 0$.
\end{definition}

\begin{proposition}[Log-space form]
\label{prop:logspace}
With $\alpha = (q - a_j)/\sigma_j$ and $\beta = (p - a_j)/\sigma_j$,
\begin{equation}
\log S_j(p) = -\operatorname{softplus}(\beta),
\qquad
\frac{S_j(q)}{S_j(p)} = \exp\!\big(\operatorname{softplus}(\beta) - \operatorname{softplus}(\alpha)\big).
\label{eq:logspaceratio}
\end{equation}
\end{proposition}

\begin{proof}
$S_j(p) = 1/(1+e^{\beta})$ by \eqref{eq:S}, so $\log S_j(p) = -\log(1+e^{\beta})$, which
is the first claim. The ratio is the exponential of the difference of the logs.
\end{proof}

The implementation evaluates $\operatorname{softplus}$ as \code{math.Log1p(math.Exp(x))}
below a clamp of $30$ and as $x$ above it. That substitution is safe: at $x = 30$,
$\log(1+e^{30}) - 30 = e^{-30} \approx 9.4\times10^{-14}$, a relative error under
$10^{-14}$.

\subsubsection*{The wrong way, and why it is so tempting}

The obvious defensive move is to clamp \emph{each exponent} at $\pm 30$ before
exponentiating, compute $S_j(q)$ and $S_j(p)$ separately, and divide. This overflows
nothing and looks careful. It is catastrophically wrong, and the direction of the error
is the worst possible one.

Take a genuinely tight player who has fallen: $a_j = 10$, $\sigma_j = 0.5$, and a
five-pick horizon. By \eqref{eq:tailhazard} the true conditional survival is
$e^{-5/0.5} = e^{-10} = 4.54\times10^{-5}$, and by Proposition~\ref{prop:tail} it stays
there no matter how much further he falls --- the hazard is memoryless in the tail.
Now watch what per-$S$ clamping does as he falls (Table~\ref{tab:clampbug}).

\begin{table}[htbp]
\centering
\small
\begin{tabular}{rrrrr}
\toprule
$p_{\text{now}}$ & $\beta$ & $\alpha$ & log-space \eqref{eq:logspaceratio} & clamp-each-$S$ \\
\midrule
20 & 20 & 30 & $4.54\times10^{-5}$ & $4.54\times10^{-5}$ \\
22 & 24 & 34 & $4.54\times10^{-5}$ & $2.48\times10^{-3}$ \\
25 & 30 & 40 & $4.54\times10^{-5}$ & $\mathbf{1.000}$ \\
30 & 40 & 50 & $4.54\times10^{-5}$ & $\mathbf{1.000}$ \\
\bottomrule
\end{tabular}
\caption{The clamping bug, with $a_j = 10$, $\sigma_j = 0.5$, horizon $5$ picks. The
truth is constant. Clamping each survival separately saturates both exponents at the
same clamped value, so their ratio collapses to $1$: the model reports a player who has
fallen twenty picks past his price as \emph{certain} to survive.}
\label{tab:clampbug}
\end{table}

The mechanism is simple once you see it. Clamping does not distort $S_j(q)$ and
$S_j(p)$ by similar amounts --- it distorts them to the \emph{same value}, because both
exponents get pinned to the clamp. The ratio of two identical numbers is $1$. And the
failure is monotone in exactly the wrong direction: \textbf{the further past his ADP a
player has fallen, the safer the broken model says he is.} A tool whose entire purpose
is to flag fallers would have been reporting $100\%$ survival for precisely the players
it was built to reason about.

The log-space form has no such failure mode because the difference
$\operatorname{softplus}(\beta) - \operatorname{softplus}(\alpha)$ degrades gracefully:
past the clamp both softpluses are their arguments, so the difference is
$\beta - \alpha = -(q-p)/\sigma_j$ exactly --- which by Proposition~\ref{prop:tail} is
precisely the right answer. The approximation is not merely safe, it is the correct
asymptotic.

The lesson generalises: \emph{when you need a ratio of tiny numbers, never compute the
numbers}. Compute the difference of their logarithms. A clamp applied before a
subtraction destroys the very quantity you are subtracting for.

\subsubsection*{Two more guards}

\begin{itemize}
\item \textbf{Missing ADP.} A live draft feed registers handcuffs and rookies no ADP
source ranked. Those players get $a_j = 999$ (\code{UndraftedADP}) --- off the drafted
radar, survives everything. Zero would be catastrophic: it would put them at the very
top of the board's price curve.

\item \textbf{Horizons in the past.} On a finished draft, $q_1$ falls back to the final
pick, which is behind $p_{\text{now}}$. Unguarded, $\alpha < \beta$ and the ratio
exceeds $1$; a replay frame prints ``$105\%$''. The engine clamps $q \gets
\max(q, p_{\text{now}})$, which restores the correct answer: no picks intervene, so
everyone survives with probability $1$.
\end{itemize}

\subsection{Falling players are flagged, not model-fixed}

There is a real temptation, on seeing a player twenty picks past his price, to adjust
his ADP. The engine does not, and the reason is a discipline worth stating: our ADP is
a measurement, and a measurement that disagrees with reality is information, not an
error to be smoothed away.

Instead, the engine marks him. Player $j$ is \emph{falling} when
\begin{equation}
\frac{p_{\text{now}} - a_j}{\sigma_j} \;\ge\; \texttt{FallerSigmas} = 1.0,
\label{eq:falling}
\end{equation}
i.e.\ when the draft has run a full one of \emph{his own} standard scale units past his
price. Measuring in his own $\sigma$ is the point: eight picks past ADP is seismic for a
first-rounder with $\sigma = 1$ and unremarkable for a round-13 flier with $\sigma = 20$.

The UI renders a falling player's numbers amber. His survival stays honest --- by
Proposition~\ref{prop:conditional} it is already correct, and by
Proposition~\ref{prop:tail} it is already generous --- and his urgency needs no special
case. The flag is the entire feature.

\subsection{Sigma shrinkage: empirical Bayes on the spread}
\label{sec:shrink}

\subsubsection*{The problem, with a real example}

The New York Jets defense, on the shipped 2026 board: ADP $174.8$, spread
$\varsigma = 1.8$, \code{times\_drafted} $= 5$. Equation \eqref{eq:sigmaclamp} turns that
into $\sigma = 1.8/1.8138 = 0.99$, and \eqref{eq:perpick} turns \emph{that} into a
per-pick hazard of $63\%$. The model is asserting that it knows, to within about one
pick, when a round-15 defense will be drafted.

It knows no such thing. It has five observations. Five drafts that happened to agree is
what five drafts usually look like when the thing being measured is a round-15 defense
nobody thinks about --- and the same board has players with two hundred observations
behind their spread. The two numbers should not be trusted equally, and
\eqref{eq:sigmaclamp} trusts them identically.

Across the 2026 pool, \code{times\_drafted} runs $5 / 49 / 215$ (minimum / median /
maximum) and $53$ of $203$ players are supported by fewer than $25$ drafts. Be careful
with the headline number here: FFC's snapshot covers $1{,}187$ drafts, but that is the
size of the \emph{window}, not the support behind any one player's spread --- a player
only appears in the drafts he was actually taken in, and the deepest names were taken in
five of them. The spread that matters is $5$ against $215$, a factor of $43$, and that is
the spread the pseudo-count $n_0$ has to be read against.

\subsubsection*{Shrinkage, and why it applies to the variance}

\begin{definition}[Empirical Bayes]
An estimate is \emph{shrunk} when it is pulled toward a central value; \emph{empirical
Bayes} means that central value (the \emph{prior}) is estimated from the same data
pool rather than supplied from outside. The classic result --- Stein's paradox, made
practical by Efron and Morris --- is that shrinking many noisy estimates toward a common
centre beats using each one raw, even though each raw estimate is individually
unbiased.
\end{definition}

The football reading: \emph{a spread computed from five drafts is mostly noise; one from
two hundred is signal. Blend each player's measured spread with the typical spread for
players priced where he is, weighted by how much data he actually has.}

The blend is on \textbf{variance}, not on standard deviation. Two reasons, one
statistical and one arithmetic.

\emph{Statistical.} Suppose $\varsigma_j^2$ is a sample variance from $n_j$
observations, and put a conjugate prior on the underlying variance --- a scaled
inverse-$\chi^2$ with $n_0$ degrees of freedom and scale $\tau^2$. The posterior for
the variance is again scaled inverse-$\chi^2$, with degrees of freedom $n_j + n_0$ and
scale
\begin{equation}
\hat\varsigma_j^{\,2} \;=\; \frac{n_j\,\varsigma_j^2 + n_0\,\tau^2}{n_j + n_0}.
\label{eq:shrinkvar}
\end{equation}
This is exactly a sample-size-weighted average of the observation and the prior, and it
is what conjugacy buys you: the arithmetic of ``how much should I believe my own
measurement'' comes out as a weighted mean with the weights being counts. (Being
scrupulous: writing $\nu = n_j + n_0$ for the posterior degrees of freedom, the posterior
\emph{mean} of the variance carries an extra factor $\nu/(\nu-2)$ over the posterior
\emph{scale}, and we use the scale. At $n_0 = 25$ every player on the board has
$\nu \ge 30$, so that factor runs from $1.07$ at the thinnest-sampled player down to
$1.03$ at the median, and $n_0$ is a
tuned pseudo-count rather than a derived one, so the distinction is well inside the
noise we are already living with. It is stated here so nobody mistakes
\eqref{eq:shrinkvar} for a theorem it is not.)

\emph{Arithmetic.} Variance is the quantity that adds. Standard deviations do not
combine linearly, and averaging square roots is not the square root of the average ---
by Jensen's inequality, since $\sqrt{\cdot}$ is concave, averaging the roots always
lands \emph{below} the root of the average. Concretely for the Jets defense with
$\tau = 18.84$ (see below) and $n_0 = 25$:
\[
\text{variance blend: } \sqrt{\tfrac{5(1.8)^2 + 25(18.84)^2}{30}} = 17.21,
\qquad
\text{stdev blend: } \tfrac{5(1.8) + 25(18.84)}{30} = 16.00.
\]
The variance blend is the more conservative of the two, which for a spread estimate ---
where under-stating spread means over-stating your own certainty --- is the direction
you want to err.

\subsubsection*{The prior is a line, not a constant}

What is ``the typical spread for a player priced here''? Not the pool median: spread
grows with depth, measurably. Fitting ordinary least squares of $\varsigma$ against
$a$ over each pool gives
\begin{align}
\text{2024 pool ($178$ players):}\quad &\tau(a) = 0.76 + 0.0876\,a, \quad \text{median } 6.9, \label{eq:prior24}\\
\text{2026 pool ($203$ players):}\quad &\tau(a) = 2.19 + 0.0953\,a, \quad \text{median } 10.3. \label{eq:prior26}
\end{align}
On the 2026 line, the prior expects $3.1$ picks of disagreement about the ADP-10 player
and $18.8$ about the ADP-175 player. Shrinking a thin sample toward a single pooled
median instead would tell a round-15 flier that the market agrees about him as tightly
as it does about a round-5 back --- which is the same mistake as the raw estimate,
merely in the other direction. Where the fitted line would predict a non-positive
spread (an extrapolation past the ends of the pool), the pool median stands in.

The prior must be fitted \emph{per pool}. Fitting on 2026 and scoring 2024 would grade
the shrink against a market it never saw --- the same era mismatch that makes the two
2025 drafts unusable in \S\ref{sec:backtest}.

\subsubsection*{What it does to the board}

\begin{table}[htbp]
\centering
\small
\begin{tabular}{lrrrrr}
\toprule
$\sigma$ across the 2026 board & min & p25 & median & p75 & max \\
\midrule
raw, from \eqref{eq:sigmaclamp}      & 0.50 & 3.61 & 5.68 & 8.71 & 21.61 \\
shrunk, $n_0 = 25$                   & 0.59 & 4.17 & 6.48 & 8.94 & 14.08 \\
\bottomrule
\end{tabular}
\caption{Sigma quantiles before and after shrinkage, $203$ players
(quantiles by linear interpolation). Players pinned at the $\sigma_{\min} = 0.5$ floor:
$3$ before, $0$ after. The shrink pulls in both tails, but it is not symmetric --- the
top end moves much more, because extreme spreads live on thinly-sampled deep players.}
\label{tab:shrink}
\end{table}

The Jets defense goes from $\sigma = 0.99$ to $\sigma = 9.49$, i.e.\ from a $63\%$
per-pick hazard to $10\%$. That is the correction the whole subsection exists for. Note
that the shrink does not \emph{always} widen: a well-sampled player with an unusually
large spread gets pulled in, which is why the maximum drops from $21.61$ to $14.08$.
The shrink is not a safety margin, it is a weighted average, and it moves whichever way
the evidence is thin.

\S\ref{sec:results} grades it. On the raw untilted model, shrinkage moves Brier from
$0.0868$ to $0.0856$ and log-loss from $0.3007$ to $0.2895$. Once the tilt of
\S\ref{sec:tilt} is applied on top --- which is what ships, so it is the comparison that
counts --- it moves Brier from $0.0677$ to $0.0670$ and log-loss from $0.2288$ to
$0.2250$. Small, consistent, and in the right direction on every horizon segment. It
ships not because it moved the headline but because it is right, it is cheap, and the
Jets-defense failure it fixes is a real one that a $63\%$ per-pick hazard would have
printed on screen.
%=======================================================================
\section{The exactly-\texorpdfstring{$N$}{N} tilt}
\label{sec:tilt}
%=======================================================================

\subsection{The problem: a budget the model does not respect}
\label{sec:tiltproblem}

Write $p_j = \Pr[A_j(q_1) \mid A_j(p_{\text{now}})]$ for the conditional survival of
each available player from \eqref{eq:psurvive}. Each player is either removed before my
turn or he is not, so under the model the number of removals is
$X = \sum_j \mathbf{1}[\text{$j$ is taken}]$, a sum of indicators with
$\Pr[\text{taken}] = 1 - p_j$. By linearity of expectation --- which needs no
independence assumption at all ---
\begin{equation}
\mathbb{E}[X] \;=\; \sum_{j \text{ available}} (1 - p_j).
\label{eq:expectedremovals}
\end{equation}

But we \emph{know} $X$. Exactly $N = q_1 - p_{\text{now}}$ picks happen before my turn,
and each one removes exactly one player. $X$ is not a random variable at all; it is a
constant we can read off the clock. And the model does not respect it.

How badly? Measured two independent ways:

\begin{itemize}
\item \textbf{On the live 2026 board.} Driving the scripted mock draft to completion from
all twelve seats gives $1{,}969$ board frames on which somebody is waiting. On $1{,}945$
of them \eqref{eq:expectedremovals} exceeds $N$, and averaged over frames it expects
$11.0$ removals against a mean window of $7.9$ picks. The over-prediction is not an
occasional artifact of one board state; it is what the model does almost everywhere.

\item \textbf{On the 2024 backtest} ($15{,}923$ predictions, $168$ vantages, see
\S\ref{sec:backtest}). Averaged over vantages, the target was $12.0$ picks, the model
expected $16.8$ removals, and the board actually lost $11.0$ ranked players. Mean
predicted survival was $0.8231$ against an observed rate of $0.8844$.
\end{itemize}

And critically, the error is \textbf{one-directional}. Table~\ref{tab:premiscal} is the
reliability table before any correction: it groups predictions into ten bins by what the
model said, and shows what actually happened in each bin.

\begin{table}[htbp]
\centering
\small
\begin{tabular}{lrrr}
\toprule
bin & mean predicted & observed & $n$ \\
\midrule
0.0--0.1 & 0.0274 & 0.3413 & 1383 \\
0.1--0.2 & 0.1481 & 0.6432 & 426 \\
0.2--0.3 & 0.2495 & 0.6520 & 319 \\
0.3--0.4 & 0.3475 & 0.7021 & 292 \\
0.4--0.5 & 0.4491 & 0.7190 & 306 \\
0.5--0.6 & 0.5496 & 0.7419 & 279 \\
0.6--0.7 & 0.6504 & 0.7710 & 310 \\
0.7--0.8 & 0.7527 & 0.7995 & 414 \\
0.8--0.9 & 0.8544 & 0.8678 & 643 \\
0.9--1.0 & 0.9931 & 0.9842 & 11551 \\
\midrule
overall & 0.8231 & 0.8844 & 15923 \\
\bottomrule
\end{tabular}
\caption{The original miscalibration. \emph{Every} bin below $0.9$ under-predicts
survival, several of them by more than a factor of two. This is not noise --- noise
would scatter --- it is a systematic bias, and it is the best single illustration of
what a reliability table is for.}
\label{tab:premiscal}
\end{table}

Read the top row at the table: of the $1{,}383$ times the engine said ``essentially
gone, $2.7\%$'', the player was actually still sitting there $34\%$ of the time. That is
not a model being unlucky. That is a model that is wrong in a fixable way.

\subsection{Why it happens}

The removals are not independent, and they are not free. Whoever is sitting between me
and my next turn, they make exactly $N$ picks between them, and each pick removes exactly
one player. The model treats each player's fate as a separate
coin flip against the market's price curve, and those coin flips collectively imagine a
draft in which sixteen or seventeen players vanish in twelve picks. They cannot. The
picks are a \emph{budget}, and the independence assumption spends more than the budget
has.

There is a second, smaller leak worth naming since it appears in the numbers above: the
board only holds the players ADP ranked, and a real draft spends some of its picks on
unranked names. That is why the target ($12.0$) and what the board really lost ($11.0$)
differ. We do not correct for this --- targeting $N$ rather than $11.0$ is a deliberate
choice not to leak the labels into the model --- and the residual over-correction it
implies is small and in the safe direction. The live engine has the same leak, so
correcting it here would grade a model the tool does not run.

\subsection{The fix: one scalar on the cumulative hazard}

\begin{definition}[Cumulative hazard]
For a survival probability $p_j \in (0,1]$, the \emph{cumulative hazard} over the window
is $H_j = -\log p_j \ge 0$. It is the total accumulated risk: $p_j = e^{-H_j}$, and
by \eqref{eq:tailhazard}, in the tail $H_j = (q_1 - p_{\text{now}})/\sigma_j$, i.e.\
picks-in-the-window divided by his per-pick scale.
\end{definition}

Now the correction. Multiply every player's cumulative hazard by one common factor $c$:
\[
H_j \;\longmapsto\; c\,H_j
\qquad\Longleftrightarrow\qquad
p_j \;\longmapsto\; e^{-cH_j} = p_j^{\,c}.
\]
Raising every survival probability to a common power \emph{is} scaling every cumulative
hazard by a common factor. This is the statistical model known as \emph{proportional
hazards}: it says the room drafts uniformly hungrier ($c > 1$) or uniformly slower
($c < 1$) than the market curve assumes, without changing anybody's relative position.
One knob, and it is exactly the knob the situation calls for.

Choose $c$ to make the budget balance:

\begin{definition}[The tilt]
\begin{equation}
\tilde p_j = p_j^{\,c^\star},
\qquad\text{where } c^\star \text{ solves }\quad
f(c) \;:=\; \sum_{j \text{ available}} \big(1 - p_j^{\,c}\big) \;=\; N.
\label{eq:tilt}
\end{equation}
\end{definition}

\subsection{The root exists, is unique, and bisection finds it}

\begin{theorem}[Well-posedness of the tilt]
\label{thm:tilt}
Let $p_1, \dots, p_m \in [0,1]$ and $f(c) = \sum_{j=1}^m (1 - p_j^{\,c})$ on $c \in
(0,\infty)$, with the convention $0^c = 0$ for $c>0$. Then:
\begin{enumerate}
\item[(i)] $f$ is continuous on $(0,\infty)$.
\item[(ii)] $f$ is strictly increasing provided at least one $p_j \in (0,1)$.
\item[(iii)] Writing $m_0 = \#\{j : p_j = 0\}$ and $m_1 = \#\{j : p_j < 1\}$,
\[
\lim_{c \to 0^+} f(c) = m_0, \qquad \lim_{c\to\infty} f(c) = m_1 .
\]
\item[(iv)] Hence for every $N \in (m_0, m_1)$ the equation $f(c) = N$ has exactly one
solution $c^\star \in (0,\infty)$, and bisection on any bracket $[\lo, \mathrm{hi}]$
with $f(\lo) \le N \le f(\mathrm{hi})$ converges to it.
\end{enumerate}
\end{theorem}

\begin{proof}
(i) Each term $c \mapsto 1 - p_j^{\,c} = 1 - e^{c\log p_j}$ is continuous in $c$ for
$p_j > 0$; for $p_j = 0$ the term is identically $1$ on $c>0$. A finite sum of
continuous functions is continuous.

(ii) Differentiate term by term. For $p_j \in (0,1)$,
\[
\frac{d}{dc}\big(1 - p_j^{\,c}\big) = -p_j^{\,c}\log p_j > 0,
\]
strictly, because $p_j^{\,c} > 0$ and $\log p_j < 0$. For $p_j = 1$ the term is
identically $0$ and contributes derivative $0$; for $p_j = 0$ it is identically $1$ and
also contributes $0$. So
\begin{equation}
f'(c) = -\sum_{j\,:\,0<p_j<1} p_j^{\,c}\log p_j \;>\; 0
\label{eq:tiltderiv}
\end{equation}
whenever at least one $p_j$ lies strictly between $0$ and $1$, which makes $f$ strictly
increasing.

(iii) For $p_j \in (0,1]$, $p_j^{\,c} = e^{c \log p_j} \to 1$ as $c \to 0^+$, so those
terms vanish and only the $p_j = 0$ terms survive, each contributing $1$: the limit is
$m_0$. As $c \to \infty$, $p_j^{\,c} \to 0$ for every $p_j < 1$ while $p_j = 1$ stays
$1$, so each of the $m_1$ terms with $p_j < 1$ contributes $1$: the limit is $m_1$.

(iv) By (i)--(iii), $f$ is a continuous strictly increasing function whose range is
exactly the open interval $(m_0, m_1)$. A strictly monotone function is injective, so
the preimage of $N$ is a single point; the intermediate value theorem supplies its
existence. Bisection is valid for any continuous function that changes sign across a
bracket, and monotonicity guarantees the standard bisection invariant
($f(\lo) \le N \le f(\mathrm{hi})$) is preserved at each halving.
\end{proof}

The implementation brackets $c \in [1/64,\,64]$ and halves until the bracket is narrower
than $10^{-6}$, which is about $26$ iterations --- microseconds. Each evaluation of $f$
reuses a precomputed array of $\log p_j$, so the bisection costs one \code{math.Exp} per
player per iteration rather than a full re-derivation of every survival curve.

\subsection{The properties that make it safe}

The tilt touches every number on the board, so it is worth being explicit about what it
cannot break.

\begin{proposition}[Fixed points and order]
\label{prop:tiltprops}
For any $c > 0$:
\begin{enumerate}
\item[(a)] $p_j = 1 \Rightarrow \tilde p_j = 1$, and $p_j = 0 \Rightarrow \tilde p_j = 0$.
\item[(b)] $p_i > p_j \Rightarrow \tilde p_i > \tilde p_j$ for $p_i, p_j \in (0,1]$.
\end{enumerate}
\end{proposition}

\begin{proof}
(a) $1^c = 1$ and $0^c = 0$ for every $c > 0$.
(b) $x \mapsto x^c = e^{c\log x}$ has derivative $c x^{c-1} > 0$ on $(0,\infty)$, so it
is strictly increasing there; apply it to $p_i > p_j$.
\end{proof}

Part (a) is what makes the tilt safe on the two boundary cases that matter. On \emph{my
own pick}, $N = 0$ and no picks intervene, so every $p_j = 1$ and the tilt cannot move
anybody --- the whole chain (survival $=1$, $\mathbb{E}[\text{best later}] =
v(\text{bestNow})$, urgency $= 0$) stays exact arithmetic rather than the output of a
bisection that returned $1.0000004$. The implementation short-circuits on the horizon
rather than on $N$ for exactly this reason. Undrafted-radar players with $a_j = 999$
likewise stay at $1$.

Part (b) is what makes the tilt safe on the board's \emph{ordering}. The tilt changes
what the numbers say; it cannot change who is more likely to survive than whom. Anyone
reading the board's sort order sees the same ranking before and after.

\subsection{The two clamps}

The bracket $[1/64, 64]$ can fail to contain a root, and both failures are real board
states rather than defensive padding. By Theorem~\ref{thm:tilt}(iii):

\begin{itemize}
\item $f(64) < N$: even at maximum hunger the pool cannot supply $N$ removals. This is
a \emph{thin board} --- a late endgame where almost every remaining player is
near-certain to survive, or a small test fixture being driven as a twelve-team draft.
Clamp $c = 64$. Physically: the model is doing its best and the picks will have to come
from players it has never heard of.

\item $f(1/64) > N$: even at minimum hunger the model expects more than $N$ removals.
The board is stacked with players the model has already written off (deep fallers with
$p_j$ near zero) while too few picks remain to take them all. Clamp $c = 1/64$.
\end{itemize}

On the shipped 2026 board neither clamp fires: over the $1{,}969$ mock frames of
\S\ref{sec:tiltproblem}, $c^\star$ ran $0.3812$ / $0.6551$ / $1.1215$ (min / median / max)
with \textbf{zero} clamped. On the 2024 backtest the same holds --- across all $168$
vantages $c^\star$ ran $0.2873$ / $0.4419$ / $1.0073$ with \textbf{zero} clamped. The
bracket is wide enough to be irrelevant on real data, which is the correct state for a
clamp to be in.

Note the direction: $c^\star < 1$ on $1{,}945$ of those $1{,}969$ frames, which means the
tilt is \emph{raising} survival probabilities almost everywhere. That matches
Table~\ref{tab:premiscal} exactly
--- every bin was under-predicting --- and it is a reassuring sign that the two
diagnostics, an internal budget check and an external backtest, are pointing at the same
defect.

\subsection{$N$ is supplied by the caller, not derived}

The correct removal count differs by horizon, and getting it from the snake arithmetic
rather than special-casing it at each call site is what keeps the two uses consistent.

\begin{definition}[Opponent picks before a horizon]
\begin{equation}
N(\text{at}) = \#\{\,p \in [p_{\text{now}}, \text{at}) : \mathrm{slot}(p) \ne \text{mySlot}\,\}.
\label{eq:oppbefore}
\end{equation}
\end{definition}

At $\text{at} = q_1$ every intervening pick belongs to somebody else, so
$N(q_1) = q_1 - p_{\text{now}} = N$. At $\text{at} = q_2$ exactly one intervening pick
is \emph{mine} --- the one at $q_1$ --- so $N(q_2) = q_2 - p_{\text{now}} - 1$. That
subtraction matters: if the second leg of the two-pick plan (\S\ref{sec:plan}) counted my
own pick as an opponent removal, it would price a player as at risk from a pick I am the
one making.

There is a subtlety in the backtest that is worth flagging so nobody reconciles the two
numbers wrongly. A backtest vantage \emph{stands on} my own pick and looks forward to my
next one, so its window $[\text{from}, \text{to})$ includes my own pick at
\code{from} --- I am one of the drafters removing a player, and the labels agree (a
player I take at \code{from} has draft position $<$ \code{to}, hence $y = 0$). Live, the
window opens on somebody else's pick and my next one closes it. Same rule --- every pick
inside the window removes exactly one player --- but the live window is one pick shorter
than any backtest row. That is the easy direction: shorter windows tilt less.

\subsection{One truth on screen}

The tilted probability $\tilde p_j$ is what everything downstream consumes: the expected
best survivor, the tier-hold probability, the ``safe to wait'' tag, and the survival
column in the data tab. An untilted probability never reaches the screen. This is a
deliberate discipline rather than an implementation detail --- if two panes quote
different odds on the same player, the reader has no way to tell which one the board's
ordering actually used, and the board stops being checkable.
%=======================================================================
\section{The expected best survivor}
\label{sec:ebest}
%=======================================================================

\subsection{The question}

Fix a position $P$ --- say running back. I am not going to get \emph{a specific} running
back at my next pick; I am going to get \emph{whoever is left}. So the quantity I want
is not any one player's value but
\[
\mathbb{E}\big[\max\{v_j : j \text{ at position $P$, still available at } q_1\}\big].
\]

\begin{definition}[Order statistic]
Given a finite collection of numbers, an \emph{order statistic} is one of them selected
by rank --- the largest, the second largest, and so on. Here the relevant one is the
maximum, and the collection is random because \emph{which} players are in it is random.
In football: ``the best guy actually left when my turn comes''.
\end{definition}

\subsection{The derivation}

Index the available players at position $P$ in decreasing value,
$v_1 \ge v_2 \ge \dots \ge v_m$, with tilted survival probabilities
$\tilde p_1, \dots, \tilde p_m$. Let $A_j$ be the event that $j$ survives to $q_1$, and
model the $A_j$ as independent with $\Pr[A_j] = \tilde p_j$.

\begin{theorem}[Expected best survivor]
\label{thm:ebest}
Define $M = \max\{v_j : A_j \text{ occurs}\}$, with $M = 0$ if no player survives. Then
\begin{equation}
\mathbb{E}[M] \;=\; \sum_{j=1}^{m} v_j\, \tilde p_j \prod_{i<j}\big(1 - \tilde p_i\big).
\label{eq:ebest}
\end{equation}
\end{theorem}

\begin{proof}
Consider the events
\[
B_j \;=\; A_j \cap \bigcap_{i<j} A_i^{\,c}
\qquad (j = 1,\dots,m),
\qquad
B_0 \;=\; \bigcap_{i=1}^{m} A_i^{\,c}.
\]
In words: $B_j$ is ``$j$ survives and everyone better is gone''; $B_0$ is ``the position
empties completely''.

These partition the sample space. They are disjoint --- for $j < k$, $B_j$ requires
$A_j$ and $B_k$ requires $A_j^{\,c}$ --- and they are exhaustive, since either no $A_i$
occurs (giving $B_0$) or there is a smallest index $j$ with $A_j$ occurring (giving
$B_j$).

On $B_j$ the maximum surviving value is $v_j$: every survivor has index $\ge j$ (the
better-indexed ones are gone by construction), and the list is sorted, so $v_j$ is the
largest value among survivors. Ties in value do not matter --- if $v_j = v_{j+1}$ the
maximum is the same number either way. On $B_0$, $M = 0$ by definition.

Therefore, splitting the expectation over the partition,
\[
\mathbb{E}[M] = \sum_{j=1}^{m} v_j \Pr[B_j] + 0\cdot\Pr[B_0].
\]
Independence gives $\Pr[B_j] = \tilde p_j \prod_{i<j}(1-\tilde p_i)$, which is
\eqref{eq:ebest}.
\end{proof}

Read the $j$-th term at the table: \emph{his value, times the probability he is the one
you end up choosing}. He has to survive, and everybody better has to be gone. That is
all \eqref{eq:ebest} says.

\subsection{Reading the pieces}

\subsubsection*{The residual mass is worth zero, and that is correct}

$\Pr[B_0]$ --- the position empties completely --- contributes nothing to
\eqref{eq:ebest}. That is not an omission. If every running back is gone, the value of
the best available running back is genuinely zero; there isn't one. The alternative,
substituting some notional deep-bench value, would be inventing a player.

\subsubsection*{The truncation is provably harmless}

The implementation walks the list with a running accumulator
$\mathrm{acc} = \prod_{i<j}(1-\tilde p_i)$, adding
$\mathrm{acc}\cdot\tilde p_j \cdot v_j$ and then updating $\mathrm{acc} \mathrel{*{=}}
(1 - \tilde p_j)$. It stops when $\mathrm{acc} < \varepsilon = 10^{-6}$.

\begin{proposition}[Truncation bound]
Truncating \eqref{eq:ebest} at the first $j$ where $\mathrm{acc} < \varepsilon$ omits at
most $\varepsilon\, v_1$.
\end{proposition}

\begin{proof}
$\mathrm{acc}$ is non-increasing in $j$ (each factor is in $[0,1]$), and the omitted tail
is $\sum_{k \ge j} v_k \tilde p_k \prod_{i<k}(1-\tilde p_i)$. Every $v_k \le v_1$ by the
sorting, and the remaining probabilities $\Pr[B_k]$ for $k \ge j$ sum to at most
$\mathrm{acc} < \varepsilon$ (they are disjoint events all contained in
$\bigcap_{i<j}A_i^{\,c}$). So the tail is at most $\varepsilon v_1$.
\end{proof}

On the shipped board $v_1 \approx 10^4$, so the bound is about $0.01$ value points ---
under a millionth of the numbers being compared.

The check sits at the \emph{top} of the loop rather than the bottom, and that placement
does real work. On my own pick every $\tilde p_j = 1$, so $\mathrm{acc}$ is exactly $0$
after the first term; the top-of-loop check then stops the walk after one term instead
of stepping through two hundred players multiplying by zero. It also means
$\mathbb{E}[M] = v_1$ \emph{exactly} on my own pick, since an exact zero times a value is
an exact zero. Urgency's exact zero (\S\ref{sec:urgency}) is arithmetic, not luck.

\subsection{What it replaced, and the discontinuity that removed}

Milestone 4 used a threshold rule:
\[
\mathrm{bestLater}(P) = \operatorname*{arg\,max}_{j \,:\, p_j \ge 0.5} v_j,
\]
the best player who clears a fifty-percent survival cutoff. It is intuitive and it is
what a human says out loud (``he'll probably still be there''). It also makes urgency a
\emph{discontinuous} function of the probabilities, and that is fatal for a board that
re-sorts on every pick.

Concretely. Suppose the position holds $v_1 = 900$ at $p_1 = 0.501$ and $v_2 = 600$ at
$p_2 = 0.95$. The threshold rule says $\mathrm{bestLater}$ is the $900$ player, so
urgency $= (v(\text{bestNow}) - 900)\cdot \mathrm{need}$. Now one pick happens
somewhere else on the board, nudging $p_1$ to $0.499$. The threshold rule now says
$\mathrm{bestLater}$ is the $600$ player and urgency jumps by $300 \cdot \mathrm{need}$
--- enough to re-sort the entire board --- on a change of two thousandths in one
probability.

Equation \eqref{eq:ebest} has no such cliff. It is a polynomial in the $\tilde p_j$, hence
continuous (indeed smooth) in all of them, so no single pick can re-order the board by
nudging somebody across a line. And it is the same question asked properly: the
threshold rule is a crude one-point approximation to the expectation, replacing a
distribution with its most likely-ish member.

\begin{example}[The worked anchor]
Three players, $v = (100, 90, 60)$ with $\tilde p = (0.2, 0.6, 0.99)$:
\[
\mathbb{E}[M] = 100(0.2) + 90(0.8)(0.6) + 60(0.8)(0.4)(0.99)
= 20 + 43.2 + 19.008 = 82.208 .
\]
With an open dedicated starter slot ($\mathrm{need} = 1$) and $v(\mathrm{bestNow}) = 100$,
urgency is $17.792$. This exact case is a table-driven test in the repository; if you
change \eqref{eq:ebest} it fails.
\end{example}

\subsection{Naming one player anyway}

Urgency prices the whole distribution, but a header that reads
``\code{chase $\to$ hall}'' needs a single name. The engine reports the \emph{modal}
best survivor: the $j$ maximising $\Pr[B_j] = \tilde p_j\prod_{i<j}(1-\tilde p_i)$ ---
the same per-player term \eqref{eq:ebest} sums. It answers the question a human actually
asks: \emph{who am I probably picking instead?}

It usually agrees with the old threshold answer. Where they differ, the modal answer is
the one \emph{higher} up the board, and there is a reason: a player near the top only
has to outlast the intervening picks, while a player further down also needs everyone
above him to go first. So a $49\%$ player above the line beats a $50\%$ player below it,
and when nobody clears $0.5$ at all the modal answer is the top of the position rather
than the deepest, safest name on it. That is the right failure mode; the old rule's was
not.

%=======================================================================
\section{Urgency and need}
\label{sec:urgency}
%=======================================================================

\subsection{The headline number}

\begin{definition}[Urgency]
\begin{equation}
\mathrm{urgency}(P) \;=\; \Big(v\big(\mathrm{bestNow}(P)\big) \;-\; \mathbb{E}\big[M_P\big]\Big)\cdot \mathrm{need}(P),
\label{eq:urgency}
\end{equation}
where $\mathrm{bestNow}(P)$ is the highest-value available player at $P$ and
$\mathbb{E}[M_P]$ is \eqref{eq:ebest} evaluated at $P$ with horizon $q_1$.
\end{definition}

This is the answer to the question in the title. The first factor is the value you lose
by waiting; the second is how much you care about this position at all. The board sorts
its position groups by \eqref{eq:urgency}, descending. The top group is the pick.

Three properties, all of which fall out of what has already been proved:

\begin{proposition}[Basic properties of urgency]
\label{prop:urgencyprops}
\begin{enumerate}
\item[(a)] $\mathrm{urgency}(P) \ge 0$ always.
\item[(b)] On my own pick, $\mathrm{urgency}(P) = 0$ for every position simultaneously.
Off my own pick, $\mathrm{urgency}(P) > 0$ strictly, provided $\mathrm{need}(P) > 0$ and
the position holds at least one available player with $v_1 > 0$.
\item[(c)] $\mathrm{urgency}(P)$ is continuous in every $\tilde p_j$.
\end{enumerate}
\end{proposition}

\begin{proof}
(a) By Theorem~\ref{thm:ebest}, $\mathbb{E}[M_P] = \sum_j v_j\Pr[B_j]$ with
$\sum_j \Pr[B_j] \le 1$ and every $v_j \le v_1 = v(\mathrm{bestNow})$. Hence
$\mathbb{E}[M_P] \le v_1 \sum_j \Pr[B_j] \le v_1$, so the bracket in \eqref{eq:urgency}
is non-negative; $\mathrm{need} \ge 0$.

(b) On my own pick $q_1 = p_{\text{now}}$, so by Proposition~\ref{prop:tiltprops}(a)
every $\tilde p_j = 1$; then $\Pr[B_1] = \tilde p_1 = 1$ and $\Pr[B_j] = 0$ for $j > 1$,
giving $\mathbb{E}[M_P] = v_1$ exactly and urgency $0$.

For the other direction, the tempting argument is wrong and it is worth seeing why. One
might say ``if any $\tilde p_j < 1$ then $\Pr[B_1] < 1$''. It does not follow: the
partition of Theorem~\ref{thm:ebest} gives $\Pr[B_1] = \tilde p_1$ \emph{exactly} --- the
product over $i < 1$ is empty --- so a small $\tilde p_j$ at some $j > 1$ says nothing at
all about $\Pr[B_1]$. The correct route goes through $B_0$ instead.

Off my own pick, $q_1 > p_{\text{now}}$, so for every player $\alpha_j > \beta_j$ in
\eqref{eq:logspaceratio}; softplus is strictly increasing, so
$\operatorname{softplus}(\beta_j) - \operatorname{softplus}(\alpha_j) < 0$ and
$p_j < 1$. The map $x \mapsto x^{c}$ fixes only $0$ and $1$, so $\tilde p_j < 1$ for
\emph{every} $j$ --- and therefore
\[
\Pr[B_0] = \prod_{i=1}^{m}\big(1 - \tilde p_i\big) \;>\; 0 .
\]
That single strict inequality is the whole argument. Since
$\sum_j \Pr[B_j] = 1 - \Pr[B_0] < 1$ and every $v_j \le v_1$,
\[
\mathbb{E}[M_P] \;\le\; v_1\big(1 - \Pr[B_0]\big) \;<\; v_1 ,
\]
so urgency is strictly positive whenever $v_1 > 0$ and $\mathrm{need}(P) > 0$.

(c) $\mathbb{E}[M_P]$ is a polynomial in the $\tilde p_j$.
\end{proof}

\begin{remark}[The exact characterisation, and a floating-point caveat]
Rearranging \eqref{eq:urgency} with $\sum_j\Pr[B_j] = 1 - \Pr[B_0]$,
\[
v_1 - \mathbb{E}[M_P] \;=\; \underbrace{v_1 \Pr[B_0]}_{\ge 0}
\;+\; \sum_{j}\underbrace{(v_1 - v_j)\Pr[B_j]}_{\ge 0},
\]
so urgency is $0$ exactly when every term vanishes: either $v_1 = 0$ (a position no
source values --- kickers and defenses before \S\ref{sec:kdef} anchors them), or all the
probability mass sits on players whose value \emph{equals} $v_1$. Part (b) is the two
cases of that identity we actually use.

The caveat is that ``$\tilde p_j < 1$'' holds in exact arithmetic, not always in
\code{float64}. A player with $a_j = \texttt{UndraftedADP} = 999$ early in a draft has
both softpluses underflowing to a difference below $2^{-53}$, so \code{PSurviveAt} returns
exactly $1.0$ and the tilt leaves it there. If such a player were the top of his position,
that position would read urgency $0$ off my own pick --- which is the right answer
(nobody is going to take him) but arrives by rounding rather than by the proof.
\end{remark}

Part (b) is the \emph{wait signal}, and it costs no extra machinery: zero urgency means
the best player at the position will keep, so there is nothing to lose by waiting. On my
own pick every urgency is zero simultaneously --- nobody can be taken in zero picks ---
which is why the board needs a tie-break for exactly that frame (\S\ref{sec:vor}).

\subsection{Need}

$\mathrm{need}(P) \in [0,1]$ is how much my roster wants another player at position $P$
right now. The league's starting lineup is a list of \emph{slots}: by default
QB / RB / RB / WR / WR / TE / FLEX / K / DEF, plus six bench spots, but the real value is
read from Sleeper's league metadata (this user's 2025 league runs \emph{two} flex slots,
which the default does not).

\begin{definition}[Need weights]
Write $R_{\text{rem}} = R - r(p_{\text{now}}) + 1$ for the number of rounds still to be
played, including the one in progress. Then
\begin{equation}
\mathrm{need}(P) =
\begin{cases}
0     & \text{if $P \in \{$K, DEF$\}$ and $R_{\text{rem}} > 3$} \quad \text{(suppression)},\\
1.0   & \text{if a dedicated $P$ slot is unfilled},\\
0.6   & \text{else if a FLEX slot is unfilled and $P$ is flex-eligible},\\
0.25  & \text{otherwise (a bench pick)}.
\end{cases}
\label{eq:need}
\end{equation}
\end{definition}

The three weights are a judgement call, not a measurement, and they should be read as
one. $1.0 / 0.6 / 0.25$ encodes ``a starter is worth about four bench players and about
one and two-thirds flex fills''. Nothing in the backtest grades them --- the backtest
grades \emph{survival}, and need does not enter survival at all --- so they are tuned by
eye against how the board reads at a real draft. If they are wrong, the tool
mis-\emph{orders} positions without mis-\emph{stating} any probability, which is the
better of the two failure modes and is why the calibration work went into survival
first.

\subsection{Slot filling is greedy, and greedy is optimal here}

To know whether a dedicated slot is unfilled, we have to assign my drafted players to
slots. The engine does it greedily: fill every dedicated slot first (a slot of position
$P$ takes the first unused player of position $P$), then fill flex slots from whoever is
left, then everyone else is bench. This is not merely convenient --- it is optimal.

\begin{proposition}[Greedy dedicated-first maximises filled starters]
\label{prop:greedy}
Let the lineup have $\ell_P$ dedicated slots at each position $P$ ($\ell$ for lineup; $s$
is still my seat) and $F$ flex slots accepting any position in a set $\mathcal{F}$, and
let me hold $n_P$ players at position $P$. Then the dedicated-first greedy assignment
fills
\[
\sum_P \min(\ell_P, n_P) \;+\; \min\Big(F,\ \textstyle\sum_{P\in\mathcal{F}} \max(0,\, n_P - \ell_P)\Big)
\]
slots, and no assignment fills more.
\end{proposition}

\begin{proof}
Greedy achieves that count by construction. For optimality, take any assignment and
apply the following exchange: if some player $x$ of position $P$ sits in a flex slot
while a dedicated $P$ slot is empty, move $x$ into the dedicated slot. The dedicated
slot accepts him (same position), the flex slot becomes empty, and the total number of
filled slots is unchanged. Each such exchange strictly decreases the number of
flex-seated players whose own dedicated slot is open, so the process terminates.

The resulting assignment saturates dedicated slots: for each $P$ it fills
$\min(\ell_P, n_P)$ of them (it cannot fill more --- there are only $\ell_P$ slots and
$n_P$ players --- and if it filled fewer, some $P$ player would be unassigned or
flex-seated with a dedicated slot open, which the exchange has eliminated). The remaining
players available for flex are exactly $\max(0, n_P - \ell_P)$ per position, and at most $F$ of
them can be seated. So no assignment exceeds the stated count, and greedy attains it.
\end{proof}

The order matters and the proposition says why: a flex-eligible player squatting in a
flex slot while his own position sits open costs nothing in that assignment but blocks a
different position later. The greedy rule sidesteps it entirely.

\subsection{Roster-aware strategy, for free}

This is worth noticing because it is the one place where the engine does something that
looks like strategy without anybody programming strategy. Take two elite running backs
in rounds 1 and 2. Both dedicated RB slots fill, so $\mathrm{need}(\text{RB})$ drops from
$1.0$ to $0.6$ (the flex is still open) and then to $0.25$. Meanwhile
$\mathrm{need}(\text{WR})$ is still $1.0$. Receivers sort to the top of the board on
their own. There is no archetype detection, no ``Zero RB'' mode, no build-type
classifier --- just \eqref{eq:urgency} multiplying by a weight that moved.

\subsection{K/DEF suppression is a policy, not a model}

$\mathrm{need}(\text{K}) = \mathrm{need}(\text{DEF}) = 0$ until three rounds remain.
This is worth calling out because it is the one constraint in the engine that is
\emph{not} an attempt to model reality. It is an editorial rule: \emph{nobody needs a
tool to tell them to draft a kicker in round 6, and the tool must never do so.}

The evidence that it is a policy and not a belief is that the engine carries a second,
looser constant for the same phenomenon. This league takes its first kicker in round 10
and its first defense in round 11 (medians: round 14 for both). So the opponent model of
engine v2 --- which is asking ``what will the room actually do'', a factual question ---
uses \code{OpponentKDefLastRounds} $= 6$, while our own recommendation uses
\code{KDefLastRounds} $= 3$. Two constants because they are two different questions:
what we are willing to recommend, and what the room actually does. Collapsing them would
make one of the two answers wrong.

\subsection{The endgame guard}

One more need modifier, near the end of the draft. Let $L$ be my remaining picks
(\emph{not} $R$, which is rounds) and $U$ my unfilled starting slots.

\begin{equation}
\text{bench weight} \;\mathrel{*{=}}\;
\begin{cases}
1    & U = 0 \quad \text{(lineup complete: every pick left is a bench pick)},\\
0    & L = U > 0 \quad \text{(every remaining pick must fill a starter)},\\
0.5  & L = U + 1,\ U > 0 \quad \text{(one pick of slack)},\\
1    & \text{otherwise}.
\end{cases}
\label{eq:endgame}
\end{equation}

The first line is easy to drop and \code{endgameSlack} does not: it short-circuits on
$U = 0$ before the arithmetic. Without it, a complete lineup with one pick left would hit
$L = U + 1$ and halve every remaining weight. Nothing would visibly break, because with
no open starters the halving is uniform across positions and the board's ordering is
unchanged --- which is exactly why the case has to be written down rather than trusted to
show itself.

The multiplier applies only to positions that fill \emph{none} of my open starting slots
--- the need function has already returned for anyone who fills a dedicated or a flex
slot, so reaching the bench weight \emph{is} the membership test. That ordering is what
lets this be a multiplier rather than a set-membership check, and it matters in the one
case that matters most: FLEX is a slot name, not a position, so ``is RB among my unfilled
starters'' reads false exactly when the flex slot an RB would fill is the thing still
open.

$L < U$ is deliberately \emph{not} handled: that lineup can no longer be completed, and
suppressing the rest of the board over a demand that cannot be met would bury the value
still on it. Nothing changes, and the board says nothing special.

%=======================================================================
\section{Tier-hold: cliffs as a probability}
\label{sec:tierhold}
%=======================================================================

\subsection{Tiers, briefly}

A \emph{tier} is a block of players at one position who are close enough in value that
which one you get barely matters. The boundary between tiers is where it starts to
matter a lot. Tiers come from a hand-transcribed rankings file where one covers a player,
and are derived from value gaps otherwise:
\[
\text{break a tier when}\quad \frac{\Delta v}{v_{\text{prev}}} > 0.10
\quad\text{and}\quad \Delta v \ge 0.015 \cdot v_{\text{top}},
\]
where $v_{\text{top}}$ is the value of the best player the derivation is walking --- the
best at the position when no rankings file is in play, and the best player the \emph{file
did not cover} otherwise, since the derivation only ever runs over the uncovered tail
(\code{deriveBlocks} reads \code{rest[0].Value}). The relative test finds cliffs anywhere
on the curve; the absolute floor stops the flat tail from shattering into singletons.
Tier $0$ means untiered --- kickers and defenses, which no source values --- and cliff
logic skips it entirely.

Two structural rules exist because of one failure mode, and the scope of each is worth
being exact about.

\emph{Oversize blocks are subdivided --- but only the ones a human drew.} A published
graphic once put $21$ receivers in ``tier 2''; a tier that large can never trigger a
cliff, so the board would never warn you about receivers at all. \code{subdivide} splits
any rankings-file block over \code{TierMaxSize} $= 8$ at its largest internal relative
drops. It is \emph{not} applied to derived blocks: those come out of
\code{deriveBlocks} and are appended untouched. On today's rankings-free board that shows
--- every tier is derived, and the derived blocks run QB $5/4/3/2/2/3/5$, RB
$2/11/7/5/11/15$, WR $6/2/3/5/5/6/6/7/29$, TE $6/4/5/4$. Three of those exceed eight and
one holds $29$ receivers. The gap is real and is listed in \S\ref{sec:threats};
the honest reading is that the guard was built for the failure it was built for, and the
value-gap path never got it.

\emph{No tier \emph{we} create may be a singleton.} A one-player tier is a permanent
``last one'' alarm, which is noise wearing an alarm's clothes, so \code{mergeRuntTiers}
folds a runt into its neighbour --- upward, into the block above it, except for a short
leading block which folds forward. Singletons the human drew are left alone: one player
alone in tier 1 is a statement about the position, not an artifact.

\subsection{The complement trick}

The question a cliff alarm should answer is: \emph{will there still be anybody from this
tier when I can act again?}

\begin{definition}[Tier hold]
Let $\mathcal{T}$ be the available players in the tier of $\mathrm{bestNow}(P)$. Then
\begin{equation}
p_{\mathrm{hold}}(P) \;=\; \Pr[\text{at least one of $\mathcal{T}$ survives}]
\;=\; 1 - \prod_{j\in\mathcal{T}} \big(1 - \tilde p_j\big).
\label{eq:tierhold}
\end{equation}
\end{definition}

The derivation is one line, and it is a technique worth internalising. ``At least one''
is awkward to compute directly --- you would have to sum over every nonempty subset,
inclusion--exclusion and all. Its complement, ``none of them'', is a single product
under independence. So compute the complement and subtract. That is the whole trick, and
it recurs everywhere in applied probability.

\subsection{Why a probability beats a headcount}

The original design fired a cliff alarm on a count: amber at two players left, red at
one. Counting cannot distinguish the two situations that matter most:

\begin{itemize}
\item \textbf{Three players, all about to go.} A tier of three where each player has
$\tilde p_j = 0.3$ holds with probability $1 - 0.7^3 = 0.657$; make it $\tilde p_j = 0.1$
and it holds with probability $0.271$. Count says ``three left, relax''. The probability
says the tier is evaporating.
\item \textbf{One player nobody wants.} A single deep player at $\tilde p_j = 0.95$ holds
with probability $0.95$. Count says ``last one, take him or lose the tier''. The
probability says it is a free wait.
\end{itemize}

These are opposite situations and the count gives them opposite-to-correct answers. So
cliff levels read \eqref{eq:tierhold}: red below $0.15$, amber below $0.5$. The count is
still \emph{reported}, because the copy quotes it and ``last one in tier 2'' is a claim
only the count can make --- but under $p_{\mathrm{hold}}$ red can fire with three players
left, so the last-one wording is kept only when the count really is one and otherwise
red reads ``tier 2 unlikely to hold --- holds $9\%$''.

\subsection{Two guards worth their space}

\subsubsection*{An untouched tier is never a cliff}

A cliff means a tier is \emph{emptying}, not that it happens to be small. If nobody has
been drafted out of the tier yet, no alarm fires whatever \eqref{eq:tierhold} says.
Without this rule, a rankings file's deliberate one-player tier 1 would print ``last
one, take him or lose the tier'' at pick 1.01 of an empty draft.

\subsubsection*{On the clock, the horizon moves}

This one is subtle and was found by measurement. \emph{Off} the clock, tier-hold prices
to $q_1$, the same horizon as the survival column beneath it, so the two agree and
nothing needs explaining. \emph{On} the clock, $q_1 = p_{\text{now}}$ and every survival
is $1$ by Proposition~\ref{prop:tiltprops}(a) --- so every tier would hold at exactly
$100\%$ and the alarm would go silent on the one frame where you are actually choosing.
Measured on the scripted mock: it fired on \emph{none} of the 15 on-the-clock vantages.

The fix is to price to $q_2$ when on the clock: passing is the decision being priced, so
the relevant horizon is the next time I could act on this tier, which is my pick after
this one. The two horizons genuinely differ by design, so the copy names the pick it is
pricing to --- a hold of $3\%$ printed directly above three survival cells reading
$100\%$ is two horizons one line apart with nothing marking which is which.

\subsection{Runs}

The run banner is pure counting over the pick feed and needs no probability at all.
Slide a window over the last $6$ picks; if $\ge 4$ share a position, a run is active for
that position. If two positions qualify, the one with more picks in the window wins. The
copy branches on the tier state --- ``POS run in progress --- $n$ left in tier $t$'' when
the tier still has players, ``POS run --- tier broke, no value left'' when it does not.

The banner needs no special logic to clear or to collapse: when a tier empties,
$\mathrm{bestNow}$ drops to the next tier and urgency falls on its own, because
\eqref{eq:urgency} is reading a different player. One caution the code carries: tier $0$
has two causes, and only one of them is ``no value left''. A position with nobody
available is genuinely exhausted; K and DEF are tier $0$ permanently, by design, with a
full board of players on it. Inferring ``tier broke'' from the tier number alone put
``def run --- tier broke, no value left'' above a group listing eight available
defenses, in exactly the rounds this league drafts them.

%=======================================================================
\section{Two-pick lookahead}
\label{sec:plan}
%=======================================================================

\subsection{The question urgency is too greedy to ask}

Equation \eqref{eq:urgency} optimises one pick at a time. The question at the turn is
about two: \emph{receiver now and running back on the way back, or the reverse?} A greedy
board cannot see that the receiver tier will be gone by my second pick while the running
back tier holds.

Corollary~\ref{cor:turn} is why this matters unevenly. From seat 3 of 12 the two picks at
the turn are $19$ apart and then $5$ apart; the ordering of a pair genuinely changes what
you can get. From seat 6 the gaps are $13$ and $11$ and the question is much softer.

\subsection{The pair score}

Let $q_2$ be my pick after next; if $q_2 = 0$ (the draft ends first) there is no plan.
For every \emph{ordered} pair of positions $(P, Q)$ with nonzero need --- $P = Q$
allowed, since double-tapping a position at the turn is a real strategy:

\begin{equation}
\mathrm{score}(P,Q) \;=\; \underbrace{v\big(\mathrm{bestNow}(P)\big)\cdot \mathrm{need}^{\circ}(P)}_{\text{first leg}}
\;+\; \underbrace{\mathbb{E}\big[M_Q^{\,-x}\big]\big|_{q_2}\cdot \mathrm{needAfter}(Q \mid x)}_{\text{second leg}},
\qquad x = \mathrm{bestNow}(P).
\label{eq:plan}
\end{equation}

The recommendation is the argmax over pairs. With at most six positions that is at most
$36$ pairs and no search --- but it is not free: each pair re-solves the tilt inside the
loop, so \code{BestPlan} measures at $1.3$--$2.9$\,ms on the shipped $203$-player board
against $0.33$--$0.51$\,ms for a full six-position urgency pass. A render happens on a
pick or a keypress, so that is affordable and the formula stays as written rather than
hoisting the solve out and drifting from it. Four details:

\begin{itemize}
\item $\mathrm{need}^{\circ}$ is not quite $\mathrm{need}$. Membership in the candidate
set is decided by \eqref{eq:need} in full --- the same number that decides whether the
pane below the plan line shows the position at all, so the plan can never name a group
the reader cannot see. The \emph{weight} is the slack-free version, \eqref{eq:need}
without the endgame multiplier of \eqref{eq:endgame}, written $\mathrm{need}^{\circ}$
here and \code{needFrom} in the code. The reason is symmetry: the second leg is priced by
\code{NeedAfter}, which carries no slack either, and mixing the two made the score depend
on the order of two legs that finish at the same roster.

\item $\mathrm{needAfter}(Q \mid x)$ is $\mathrm{need}^{\circ}$ recomputed as if $x$ were
already on my roster. Taking a running back with the first pick is exactly what drops RB
to flex weight for the second, and the pair score has to see that. It is computed on a
\emph{copy} of the roster --- appending onto the live slice would write into its spare
capacity, invisible to any equality check on the state and still a mutation during a
render.

\item $\mathbb{E}[M_Q^{-x}]$ excludes $x$ from the pool: I took him. A no-op unless
$P = Q$, and that is precisely the case worth getting right.

\item The second leg's tilt uses $N(q_2) = q_2 - p_{\text{now}} - 1$ from
\eqref{eq:oppbefore}, because my own pick at $q_1$ is not a removal somebody else makes.
Getting this wrong prices the second leg against a phantom opponent who is me.
\end{itemize}

\subsection{The honest simplification}

The first leg is priced at $v(\mathrm{bestNow}(P))$ --- \emph{today's} best available ---
even when my next pick is seventeen picks away and that player will plainly be gone by
then. This is a deliberate simplification, not an oversight, and it is stated here
because it is the single largest approximation in the engine.

The defence is that \eqref{eq:plan} is a \emph{ranking}, not a forecast. Every pair is
scored against the same board and carries exactly the same optimism in its first leg, so
the comparison between pairs survives even though the absolute number does not mean
anything. Which is why the score is deliberately \emph{not} rendered on screen: printed
as ``(e 12318)'' it read as an expectation while sitting three rows above the same
player's survival column reading $0\%$ --- one line of the board contradicting the
evidence directly under it, on $27$ of $166$ plan frames of the scripted mock. The pair
is the product; the number was noise wearing a label it could not honour.

A proper fix exists and is scoped as future work: search over \emph{my} picks inside the
Monte Carlo rollout that engine v2 already runs for opponents. That is milestone 7
material, not a patch.

\subsection{The one thing that outranks the score}

Late in a draft, finishing the lineup beats maximising the score, and by a lot, because
K and DEF carry no imported value --- both legs of any pair containing one score near
zero, so such a pair can never win an argmax over value. Measured: at pick 14.10 with a
defense still to fill and exactly two picks left, the plan line read ``rb at 14.10
$\to$ rb at 15.03'': two bench running backs and no defense at all.

So the plan first maximises the number of legs that fill an open starting slot, clamped
to what is arithmetically required ($\mathrm{mustFill} = \operatorname{clamp}(2 - (L - U),
0, 2)$, with $L$ and $U$ as in \eqref{eq:endgame}), and only then maximises
\eqref{eq:plan}. Up to the last couple of rounds $\mathrm{mustFill} = 0$ and this is a
plain argmax on score.
%=======================================================================
\section{Room-warped prices}
\label{sec:room}
%=======================================================================

\subsection{The idea}

National ADP is the average of thousands of strangers. A fixed set of twelve
leaguemates is not the average of anything. If this room reliably takes quarterbacks
seventeen picks earlier than the national market prices them, then every survival
probability computed from national ADP is wrong for quarterbacks in a knowable
direction.

We have this league's own completed drafts: three of them, $552$ picks. Crucially, the
curve we build from them reads \emph{only pick order and position} --- never a player
name, never a historical ADP. So the era caveat that makes the two 2025 drafts unusable
for scoring survival (\S\ref{sec:backtest}) does not apply here. ``The fourth
quarterback went at pick $34.7$'' is as true in 2024 as in 2026: it is a statement about
the room's behaviour, not about any player.

\subsection{The room's price curve}

\begin{definition}[Room curve]
Write $\mu(P, k)$ for the mean, over the drafts that reached that depth, of the overall
pick at which the $k$-th player of position $P$ was taken. Then
\begin{equation}
\adp_{\text{room}}(P, k) \;=\; \max_{k' \le k}\; \mu(P, k').
\label{eq:roomcurve}
\end{equation}
The running max is over $k'$ and the quantity inside depends on $k'$ --- that is the
whole content of the definition, and \S\ref{sec:roommono} is why it is there.
\end{definition}

Table~\ref{tab:roomcurve} is the curve over all three drafts.

\begin{table}[htbp]
\centering
\small
\begin{tabular}{lrrrrrrrrr}
\toprule
pos & depth & $k{=}1$ & 2 & 3 & 4 & 5 & 6 & 7 & 8 \\
\midrule
QB  & 21 & 23.0 & 27.0 & 31.7 & 34.7 & 39.0 & 58.0 & 68.0 & 81.3 \\
RB  & 59 & 1.7 & 3.3 & 4.3 & 7.0 & 9.0 & 11.3 & 13.7 & 15.3 \\
WR  & 70 & 1.3 & 4.3 & 6.3 & 8.7 & 10.0 & 11.0 & 13.3 & 15.3 \\
TE  & 19 & 23.3 & 31.7 & 41.0 & 56.3 & 63.7 & 70.0 & 75.0 & 78.3 \\
K   & 12 & 117.7 & 123.3 & 132.3 & 136.7 & 144.7 & 153.0 & 160.3 & 167.7 \\
DEF & 13 & 129.3 & 135.7 & 139.3 & 143.7 & 154.0 & 160.0 & 161.7 & 168.0 \\
\bottomrule
\end{tabular}
\caption{$\adp_{\text{room}}(P,k)$ over three completed drafts of this league
($552$ picks). \emph{depth} is the deepest $k$ the room ever reached at the position;
past it the curve has no opinion at all.}
\label{tab:roomcurve}
\end{table}

Against the 2026 national board, the mean gap over the first five at each position is:
QB $-17.4$, TE $-13.4$, K $-8.6$, WR $-0.8$, RB $+0.4$, DEF $+30.0$ picks (negative
means the room is earlier than the market). This room is emphatically quarterback- and
tight-end-early at the top of those positions, and takes defenses much later than the
market does.

\subsubsection*{The running maxes are not cosmetic}
\label{sec:roommono}

Within a single draft the $k$-th player at a position obviously cannot go before the
$(k{-}1)$-th. But the \emph{means} can invert, because a deep $k$ is supported by fewer
drafts than a shallow one: if two of three drafts stopped at 19 quarterbacks, then
$\adp_{\text{room}}(\text{QB}, 20)$ is one draft's number and can easily sit earlier
than the three-draft average in front of it. An inverted price curve would tell the
board that the 20th quarterback comes off before the 19th, which is not a small error
--- it is a statement that cannot be true.

There are in fact \textbf{three} monotonicity guards in the pipeline, each catching a
different inversion, and each was added after measuring one:

\begin{enumerate}
\item The running max inside \eqref{eq:roomcurve}, above.
\item A second running max over the \emph{blend} (below). The two inputs are each
non-decreasing in $k$, but the weight between them is not constant --- a thinly supported
$k$ stays nearer its market price than the well-supported $k$ in front of it. Measured on
the 2026 board: the 13th defense (one draft deep) priced $5.2$ picks \emph{earlier} than
the 12th (three drafts deep).
\item A \emph{ceiling}. Since only the top $K$ players at a position are warped
(\S\ref{sec:roomfail}), a warped player must not be priced later than the market prices
the next, unwarped man --- or the two halves of the sequence do not join up. Measured:
at $K = 5$ the warp priced the 5th defense at $144.16$ while the 6th kept his raw ADP of
$129.90$, so the board read the 6th defense as coming off $14.3$ picks \emph{before} the
5th.
\end{enumerate}

The general lesson: any time you splice two monotone sequences with a varying weight, or
truncate a transformation partway down a sorted list, check monotonicity of the
\emph{output}. Monotone inputs do not give it to you.

\subsection{Blending, with a pseudo-count}

We do not want to replace the market's price with three nights of evidence. We want to
move toward it in proportion to how much evidence there is.

\begin{definition}[Warped effective ADP]
For player $j$ whose position-ADP-rank among board players is $k$,
\begin{equation}
a^{\text{eff}}_j \;=\; w\cdot\adp_{\text{room}}(P,k) \;+\; (1-w)\cdot a_j,
\qquad
w = \frac{n_k}{n_k + 2},
\label{eq:warp}
\end{equation}
where $n_k$ is the number of drafts that actually reached depth $k$ at that position.
\end{definition}

\begin{definition}[Pseudo-count]
A \emph{pseudo-count} is a quantity of imaginary prior evidence added to a real count so
that a small sample cannot dominate. Here $\texttt{RoomWarpPseudo} = 2$ says ``the
national market is worth two of this room's drafts''. So three drafts buy $60\%$ of the
warp, one draft buys $33\%$, and a depth no draft ever reached buys none --- the rule
``$w \to 0$ outside the observed range'' falls out of the same formula rather than being
a special case.
\end{definition}

Two is deliberately timid. Three drafts of one room is thin evidence and the whole point
of a pseudo-count is that the answer degrades toward the market rather than toward
whatever three nights produced.

The warp is a \textbf{location shift only}. $\sigma_j$ is never touched. The room
disagreeing about a player's price is not the same claim as the room being more
predictable about him, and only one of those two is measurable from three drafts.

\subsection{Wiring: exactly two readers}

Raw ADP is read in six places in the codebase and most of them must \emph{not} warp. So
the warped number lives in a separate field \code{Player.ADPEff} rather than overwriting
\code{Player.ADP}, and precisely two functions read it: the survival curve and the
falling flag. Those are the two asking ``when does this player come off the board'',
which is exactly the question a room disagrees with the market about. The display
columns, the board's sort tie-breaks, and --- critically --- the mock draft's simulated
picker all keep the raw market price. Warping the mock's picker would warp the fake room
itself, and the demo would confirm its own hypothesis.

\subsection{Why the full-depth warp fails}
\label{sec:roomfail}

This is the most interesting negative result in the project and it deserves the space.

\subsubsection*{The structural argument}

National ADP ranks far more players at a position than any finite draft takes. The 2024
board lists $23$ quarterbacks. A 12-team, 15-round draft takes about $19$. So:

\begin{itemize}
\item Past the room's appetite, $\adp_{\text{room}}(P,\cdot)$ has \emph{no data at all}
--- the depth column of Table~\ref{tab:roomcurve} is where it stops.
\item Well \emph{before} that point it is already unreliable, because the deep $k$ values
are supported by one or two drafts rather than three.
\item And there are \textbf{more players in that unreliable tail than at the reliable
head}. Five quarterbacks at the top against eighteen behind them.
\end{itemize}

So the tail swamps the head, and it does so with a systematic bias rather than noise:
mapping rank $\to$ room-pick necessarily prices deep players \emph{later} than the market
does, because the market keeps ranking players the room simply never gets to.

\subsubsection*{The measurement}

Table~\ref{tab:qbwarp} is the held-out curve (built from the two 2025 drafts only) laid
against the 2024 national board's quarterbacks.

\begin{table}[htbp]
\centering
\small
\begin{tabular}{llrrr}
\toprule
rank & player & national ADP & room pick & room $-$ national \\
\midrule
QB1  & Josh Allen             & 27.3 & 24.0 & $-3.3$ \\
QB2  & Jalen Hurts            & 33.9 & 26.5 & $-7.4$ \\
QB3  & Patrick Mahomes        & 36.2 & 31.5 & $-4.7$ \\
QB4  & Lamar Jackson          & 40.6 & 34.5 & $-6.1$ \\
QB5  & C.J. Stroud            & 44.2 & 40.0 & $-4.2$ \\
\midrule
QB6  & Dak Prescott           & 56.7 & 63.0 & $+6.3$ \\
QB8  & Anthony Richardson Sr. & 64.4 & 84.5 & $+20.1$ \\
QB10 & Jordan Love            & 78.7 & 100.0 & $+21.3$ \\
QB12 & Tua Tagovailoa         & 89.6 & 126.5 & $+36.9$ \\
QB14 & Justin Herbert         & 103.6 & 138.5 & $+34.9$ \\
\bottomrule
\end{tabular}
\caption{The sign flip. The room is genuinely earlier than the market on the top five
quarterbacks and dramatically later on everyone behind them --- because the market keeps
ranking quarterbacks the room never reaches. $23$ quarterbacks on the board; about $19$
taken per draft.}
\label{tab:qbwarp}
\end{table}

\subsubsection*{The consequence, which is the opposite of the intent}

Table~\ref{tab:qbwarp} shows the raw curve; what the model actually sees is that gap
multiplied by the blend weight $w \approx 0.5$ of \eqref{eq:warp}. Averaged over every
prediction, the uncapped warp shifts named quarterbacks \textbf{$+11.0$ picks later}. The warp was built to capture ``this room is
quarterback-early''. Applied at full depth, it says quarterbacks go \emph{later} than the
market thinks --- the exact opposite of the signal it exists for --- because the eighteen
tail quarterbacks outvote the five at the head.

And it loses the gate: Brier $0.0670 \to 0.0671$, log-loss $0.2250 \to 0.2327$. Both
worse.

\subsubsection*{Rank versus player: the apparent contradiction, resolved}

There is something that looks like a contradiction in the measurements and it is worth
untangling, because it is the sharpest thing the data says about this room.

\begin{itemize}
\item \textbf{QB \emph{slots} go earlier than national.} The room takes its QB1 at pick
$23.0$ against a market ADP of $26.1$, and the top five run $17.4$ picks early.
\item \textbf{QB \emph{players} are the worst-calibrated position} in the 2024 backtest.
Under the shrunk-plus-tilt model the warp is gated against --- the numbers
\code{calibrate} prints, and the same pair quoted in \S\ref{sec:results} --- quarterbacks
are predicted at $0.841$ against an observed $0.917$ --- an under-prediction of $0.076$,
four times the next-largest ($+0.019$, running backs). Named quarterbacks survive
\emph{longer} than the model expects.
\end{itemize}

Both are true, and they are not in tension once you separate the rank from the man. The
room reaches for quarterbacks early --- but not for the \emph{same} quarterbacks the
national board ranks first. This room's QB4 is not the market's QB4. So the position's
slots empty early while any particular named quarterback lasts longer than his national
price implies, because the reaches are scattered across a wider set of names.

A rank-based warp captures the first fact. It cannot capture the second, and at depth it
actively fights it.

\subsubsection*{The fix, and its sweep}

Restrict the warp to the top $\texttt{RoomWarpTopK} = 5$ players at each position;
everyone deeper keeps raw national ADP. Then it passes, and comfortably: Brier
$0.0670 \to 0.0660$, log-loss $0.2250 \to 0.2222$. Sweeping the cutoff on the same
backtest:

\begin{center}
\small
\begin{tabular}{lrrrrrr}
\toprule
cutoff $K$ & 1 & 3 & 5 & 8 & 12 & uncapped \\
\midrule
Brier    & 0.0665 & 0.0662 & \textbf{0.0660} & 0.0661 & 0.0674 & 0.0671 \\
log-loss & 0.2237 & 0.2227 & \textbf{0.2222} & 0.2229 & 0.2283 & 0.2327 \\
\bottomrule
\end{tabular}
\end{center}

The optimum is a plateau at $K = 4$--$5$ and the sign flip is well past it, which is what
you want to see: a knife-edge optimum would suggest a fitted artifact, a plateau suggests
a real effect with a real boundary.

Per position at $K = 5$: five of six improve on both metrics outright; quarterbacks
improve on both (Brier $0.0970 \to 0.0952$, log-loss $0.2966 \to 0.2921$), and receivers
improve on Brier while their log-loss moves $+0.000032$, which prints as
$0.2450 \to 0.2450$. That is a tie the reader can check in the output, not a win, and the
tool's own gate scores it as a tie and says so on every run rather than leaving it to a
comment.

\subsubsection*{What would overturn it}

\textbf{The evidence is one fold.} 2024 is the only season with era-correct ADP, so two
drafts build the curve, one scores it, and the whole margin is $-0.0009$ Brier and
$-0.0028$ log-loss. It ships because it is the best number available and because no
position and no nearby cutoff points the other way --- not because a single fold can
tell a real edge from a lucky one.

A second scorable season in FFC's archive would settle it. As of this writing
\code{year=2025} still answers ``No ADP data found''. Recheck once, cheaply, near draft
day; re-gate then. If $K \le 5$ stops winning, \code{-room=false} is already the
fallback and the default moves back.

%=======================================================================
\section{Value over replacement, and pricing the unpriced}
\label{sec:vor}
%=======================================================================

\subsection{Why raw value is the wrong tie-break}

By Proposition~\ref{prop:urgencyprops}(b), every position's urgency is exactly zero on my
own pick. That is the frame where I am actually choosing, so the tie-break decides what
the board points at when it matters most.

The obvious tie-break is $v(\mathrm{bestNow}(P)) \cdot \mathrm{need}(P)$: rank each
position by its best available player. But that answers ``who is the best player left'',
which is not the question. The best receiver on a board $69$ receivers deep and the best
tight end on a board $19$ deep can carry the same value while buying completely different
things, because behind the receiver sit sixty more and behind the tight end sit two.

\begin{definition}[Value over replacement]
Let $\mathrm{dem}(P)$ be the \emph{demand} for position $P$: how many players at it this
league drafts in one draft. Let $v_{\mathrm{rep}}(P)$ be the value of the
$\mathrm{dem}(P)$-th best $P$ on the \emph{pre-draft} board. Then
\begin{equation}
\mathrm{vor}(j) \;=\; \max\big(0,\; v_j - v_{\mathrm{rep}}(\pos(j))\big).
\label{eq:vor}
\end{equation}
\end{definition}

$v_{\mathrm{rep}}(P)$ is the player you would have ended up with at that position anyway.
Subtracting it measures what taking $j$ actually \emph{buys}.

\subsection{Demand is measured, not assumed}

$\mathrm{dem}(P)$ is the median count of the position taken per draft over this league's
own cached drafts (median, not mean: one of the three drafts is 16 rounds and two are 15,
so a mean would quietly weight the long draft's extra twelve picks into every position).

\begin{table}[htbp]
\centering
\small
\begin{tabular}{lrrrl}
\toprule
pos & $\mathrm{dem}(P)$ & valued on board & $v_{\mathrm{rep}}(P)$ & the player you'd have settled for \\
\midrule
QB  & 19 & 24 & 314 & Tyler Shough (ADP 140.6) \\
RB  & 58 & 51 & 0   & nobody --- only 51 are valued, so replacement is free \\
WR  & 67 & 69 & 10  & Malik Washington (ADP 164.2) \\
TE  & 17 & 19 & 182 & Hunter Henry (ADP 146.3) \\
K   & 12 & 15 & 112 & Evan McPherson (ADP 158.2) \\
DEF & 12 & 19 & 148 & Dallas Cowboys (ADP 153.2) \\
\bottomrule
\end{tabular}
\caption{Replacement levels on the shipped 2026 board. The spread \emph{is} the content:
quarterbacks carry a $314$-point discount and tight ends $182$, against none at all for
running backs.}
\label{tab:vor}
\end{table}

Compare the fallback, which is the league's lineup \emph{shape} rather than a measurement
(every team's dedicated slots plus an equal share of each eligible flex): that gives
QB 12, RB 28, WR 28, TE 16. The gap is the bench. A league starting one quarterback
still drafts nineteen of them, and it hoards twice as many running backs as it can start.
The fallback is a floor, not an estimate, and the tool prints which of the two is in play.

Running back is the instructive row: the room drafts $58$ backs while only $51$ carry a
value at all, so replacement at running back is off the bottom of the board and a back's
vor is his whole value. That is not a degenerate case, it is what ``deep position'' means
expressed as a number.

$v_{\mathrm{rep}}(P)$ is deliberately \emph{static} --- computed over the full pre-draft board, ignoring
who has been taken. A replacement level that climbed as a position emptied would make
every position look scarcer the longer you waited, which is exactly backwards.

Exactly one call site changed when vor landed: the zero-urgency tie-break in the board's
group sort. Everything else still ranks by urgency.

\subsection{Pricing kickers and defenses}
\label{sec:kdef}

FantasyCalc values QB/RB/WR/TE only. Kickers and defenses arrive with $v_j = 0$, and a
position stuck at zero can never produce urgency however badly you need a kicker in round
15. So they need a value --- and inventing one is exactly what this project does not do.

\subsubsection*{The rejected rule: nearest ADP}

The specification said: take the value of the skill player with the nearest overall ADP,
interpolating between neighbours. Implemented literally, it produces nonsense, and the
reason is a genuinely useful observation: \textbf{value is not monotone in ADP}. Two
players priced within a fifth of a pick of each other can carry wildly different values.
Measured on today's board: Tre Tucker at ADP $128.0$ carries $9$ while Mark Andrews at
ADP $128.2$ carries $412$ --- forty-five times as much, two-tenths of a pick later. Higher
up it is the same story with bigger numbers: Trey McBride at $39.5$ carries $6{,}601$
against Davante Adams at $39.6$ with $2{,}697$. ADP measures how early the market takes a
player; value measures how good it thinks he is; scarcity at his position separates the
two, and nothing forces them to agree locally.

Interpolating between whoever happens to sit either side therefore inherits that
non-monotonicity and amplifies it. Measured on the board at the time: Denver DEF at ADP
$103.3$ got value $1698.6$, Houston DEF at $106.4$ got $989.9$, and the LA Rams at
$108.3$ got $251.1$ --- three defenses ordered nonsensically across five picks of ADP,
which then corrupts the available-player sort, since that sorts by value.

\subsubsection*{The rule that ships: anchor by rank}

Let $k$ be the number of valued skill players priced ahead of him --- that is, with
smaller ADP. Set his value to the $k$-th largest skill value.

\begin{proposition}[Monotone by construction]
The rank anchor is non-decreasing in ADP: if $a_i \le a_j$ for two K/DEF players then
their anchored values satisfy $v_i \ge v_j$.
\end{proposition}

\begin{proof}
$k$ counts skill players with smaller ADP, so $a_i \le a_j$ implies $k_i \le k_j$. The
sorted skill-value list is non-increasing, so the $k_i$-th entry is $\ge$ the $k_j$-th.
\end{proof}

It also lands \emph{exactly on} the imported curve rather than between two points of it,
and it needs no interpolation at all. Across all $34$ kickers and defenses on the shipped
board: zero monotonicity violations. Sample of the resulting curve (computed on the
current board):
\[
\text{ADP } 10 \to 8059, \quad 50 \to 3208, \quad 100 \to 868, \quad 150 \to 168, \quad 185 \to 9.
\]

\subsubsection*{The other rejected alternative, and a scale lesson}

The engine carries a rank-to-value fallback for boards with no imported values at all:
$v = \texttt{ValueBase}\cdot e^{-\mathrm{rank}/\texttt{ValueDecay}} = 250\,e^{-\mathrm{rank}/40}$.
Using it for K/DEF would have been the obvious move. It gives $250 e^{-3.75} = 5.88$ at
rank $150$, on a board where the $150$th-best skill player carries $141$ --- a
\textbf{24-fold} scale mismatch.

This is the concrete face of the project's ``never mix value modes in one draft'' rule.
Values are compared only by difference, so their scale is arbitrary --- but only if it is
\emph{one} scale. Two scales in one board is not a cosmetic problem: it would make kicker
urgency permanently invisible, since a $6$-point drop cannot compete with a $300$-point
drop no matter how much you need a kicker.

Kickers and defenses stay at tier $0$ on purpose even now that they have values. Their
value is borrowed from a skill player at the same price, so a ``gap'' between two kickers
is a gap between two receivers somewhere else on the board; it cannot draw a real tier
boundary. Tier $0$ is also what keeps cliff logic skipping them, and a red ``last one in
tier 3'' about kickers is exactly the alarm this tool must never raise.
%=======================================================================
\section{Calibration: how we know any of this}
\label{sec:calibration}
%=======================================================================

This is the intellectual core of the project. Everything above is a model; this section
is how the model is graded, and it is the reason two of the ideas in the previous
sections ship and two do not.

\subsection{Scoring a probability}

Suppose the engine says a player survives with probability $\hat p$, and then he either
does ($y = 1$) or does not ($y = 0$). (Throughout this section $\hat p$ is a
\emph{forecast} and $y$ its outcome; the letters $p$ and $q$ went to pick numbers back in
\S\ref{sec:notation} and are not reused here.) How do you score that? You cannot say ``it
was right'' or ``it was wrong'' --- a $70\%$ forecast that fails is not wrong, it is a
$30\%$ event happening.

\begin{definition}[Scoring rule]
A \emph{scoring rule} is a function $\mathcal{S}(\hat p, y)$ giving a penalty (lower is
better) for reporting probability $\hat p$ when the outcome is $y \in \{0,1\}$.
\end{definition}

\begin{definition}[Proper scoring rule]
A scoring rule is \emph{proper} if, whenever the forecaster's honest belief is that the
outcome will be $1$ with probability $\pi$, reporting $\hat p = \pi$ minimises the
\emph{expected} score
\[
\bar{\mathcal{S}}(\hat p;\pi) \;=\; \pi\, \mathcal{S}(\hat p,1) + (1-\pi)\, \mathcal{S}(\hat p,0).
\]
It is \emph{strictly proper} if $\hat p = \pi$ is the \emph{unique} minimiser.
\end{definition}

Propriety is the property that makes a metric trustworthy. Under an improper rule, a
forecaster can score better by lying --- by shading probabilities toward $0$ and $1$ to
look decisive, say --- and then the metric is measuring bravado rather than knowledge.
Under a proper rule, the only way to score well is to actually believe what you say.

\begin{definition}[Brier score]
$\mathcal{S}_{\mathrm{B}}(\hat p,y) = (\hat p - y)^2$, averaged over predictions
\cite{brier1950}.
\end{definition}

\begin{definition}[Logarithmic score / log-loss]
$\mathcal{S}_{\mathrm{L}}(\hat p,y) = -\big(y \log \hat p + (1-y)\log(1-\hat p)\big)$,
averaged over predictions, with $\hat p$ clamped to $[10^{-6}, 1-10^{-6}]$ so a single
overconfident miss does not become the entire score.
\end{definition}

\begin{theorem}[Both are strictly proper]
\label{thm:proper}
$\mathcal{S}_{\mathrm{B}}$ and $\mathcal{S}_{\mathrm{L}}$ are strictly proper.
\end{theorem}

\begin{proof}
\emph{Brier.} Expand the expected score:
\[
\bar{\mathcal{S}}_{\mathrm{B}}(\hat p;\pi) = \pi(\hat p-1)^2 + (1-\pi)\hat p^2
 = \pi \hat p^2 - 2\pi \hat p + \pi + \hat p^2 - \pi \hat p^2
 = \hat p^2 - 2\pi \hat p + \pi .
\]
This is a quadratic in $\hat p$ with positive leading coefficient. Differentiate:
$\partial_{\hat p} \bar{\mathcal{S}}_{\mathrm{B}} = 2\hat p - 2\pi$, which vanishes only
at $\hat p = \pi$, and $\partial^2_{\hat p} \bar{\mathcal{S}}_{\mathrm{B}} = 2 > 0$, so
that point is the unique minimum. Note the minimum value is $\pi - \pi^2 = \pi(1-\pi)$,
which we will need in a moment.

\emph{Logarithmic.}
$\bar{\mathcal{S}}_{\mathrm{L}}(\hat p;\pi) = -\pi\log \hat p - (1-\pi)\log(1-\hat p)$ on
$\hat p \in (0,1)$. Differentiate:
\[
\partial_{\hat p} \bar{\mathcal{S}}_{\mathrm{L}} = -\frac{\pi}{\hat p} + \frac{1-\pi}{1-\hat p},
\]
which vanishes when $\pi(1-\hat p) = \hat p(1-\pi)$, i.e.\ $\pi - \pi \hat p = \hat p -
\pi \hat p$, i.e.\ $\hat p = \pi$. The second derivative
$\pi/\hat p^2 + (1-\pi)/(1-\hat p)^2 > 0$ everywhere on $(0,1)$, so
$\bar{\mathcal{S}}_{\mathrm{L}}$ is strictly convex and the stationary point is the unique
minimum.
\end{proof}

Two proper rules rather than one, because they punish differently. Brier is bounded and
quadratic: it cares about the average size of the error. Log-loss is unbounded and blows
up as $\hat p \to 0$ on an outcome that happens: it savages confident mistakes. A change that
improves one and wrecks the other is a change that traded honest uncertainty for
false confidence, and reporting both catches it. In this project the two agreed on every
decision that mattered, which is reassuring but not guaranteed.

\subsubsection*{The coin-flip ceiling}

If you predict $\hat p = 0.5$ for everything, then $(\hat p-y)^2 = 0.25$ whatever happens,
so your Brier score is exactly $0.25$. That is why $0.25$ is quoted as the ``coin-flip
ceiling'': it is what total agnosticism buys, and anything above it is worse than saying
nothing.

But $0.25$ is a weak bar when the base rate is lopsided, and ours is. Write $\bar o$ for
the observed base rate. Predicting $\bar o$ for everything scores $\bar o(1-\bar o)$ ---
by the Brier computation above with $\pi = \bar o$. Our observed survival rate is
$0.8844$, so the honest floor is
\[
0.8844 \times 0.1156 = 0.1022,
\]
which is exactly the ``baseline: constant $0.884$'' row in the results. \emph{That} is
the number a model has to beat, not $0.25$.

\subsection{The Murphy decomposition, and what a reliability table is for}

\begin{definition}[Reliability table]
Sort all predictions into $B$ bins by the value predicted. For each bin $b$ record
$n_b$ (count), $\bar p_b$ (mean forecast in the bin) and $\bar o_b$ (observed rate in
the bin). A perfectly calibrated model has $\bar p_b = \bar o_b$ in every bin.
\end{definition}

\begin{theorem}[Murphy decomposition \cite{murphy1973}]
\label{thm:murphy}
If forecasts take finitely many distinct values and the bins group identical forecasts,
then with $n = \sum_b n_b$ the total number of predictions (\emph{not} $N$, which is
picks-until-mine) and $\bar o = \frac{1}{n}\sum_b n_b \bar o_b$ the overall base rate,
\begin{equation}
\underbrace{\frac{1}{n}\sum_{i}(\hat p_i - y_i)^2}_{\text{Brier}}
\;=\;
\underbrace{\frac{1}{n}\sum_b n_b(\bar p_b - \bar o_b)^2}_{\text{reliability}}
\;-\;
\underbrace{\frac{1}{n}\sum_b n_b(\bar o_b - \bar o)^2}_{\text{resolution}}
\;+\;
\underbrace{\bar o(1-\bar o)}_{\text{uncertainty}} .
\label{eq:murphy}
\end{equation}
\end{theorem}

\begin{proof}[Proof sketch]
Within a bin the forecast is the constant $\bar p_b$, so the bin's mean squared error is
$\frac{1}{n_b}\sum_{i \in b}(\bar p_b - y_i)^2$. Add and subtract $\bar o_b$ inside the
square and expand; the cross term is $-2(\bar p_b - \bar o_b)\sum_{i\in b}(y_i - \bar o_b)
= 0$ because $\bar o_b$ is the bin mean of $y$. That leaves
$(\bar p_b - \bar o_b)^2 + \bar o_b(1-\bar o_b)$ per bin, using $\frac{1}{n_b}\sum_{i \in
b}(y_i - \bar o_b)^2 = \bar o_b(1-\bar o_b)$ for binary $y$. Averaging over bins gives
reliability plus $\frac{1}{n}\sum_b n_b \bar o_b (1 - \bar o_b)$, and the same
add-subtract trick on the latter around $\bar o$ turns it into
$\bar o(1-\bar o) - \frac{1}{n}\sum_b n_b(\bar o_b - \bar o)^2$.
\end{proof}

Read the three terms:

\begin{itemize}
\item \textbf{Uncertainty} $= \bar o(1-\bar o)$ is a property of the \emph{world}, not of
the model. Nothing you do changes it. It is $0.1022$ here, and it is the constant-
predictor baseline.
\item \textbf{Reliability} (lower is better) is how far the model's stated probabilities
are from the truth \emph{conditional on what it said}. This is exactly the gap between
the two columns of a reliability table. A model with reliability $0$ is honest.
\item \textbf{Resolution} (higher is better, hence the minus sign) is how much the
model's bins \emph{separate} outcomes from the base rate. A model that says $0.8844$ for
everything has zero reliability error \emph{and} zero resolution: perfectly honest,
perfectly useless. Resolution is discrimination --- the ability to tell cases apart.
\end{itemize}

So Table~\ref{tab:premiscal} is a picture of a large reliability term. Every bin under
$0.9$ under-predicting means every $(\bar p_b - \bar o_b)^2$ is large and all of the
same sign: a systematic bias, which is exactly the kind of error a single scalar can fix.
That is the argument for the tilt, made visible.

Computing \eqref{eq:murphy} on our two headline models (using the ten printed bins; the
decomposition is exact only for constant-within-bin forecasts, so the reconstruction is
approximate --- it recovers $0.0879$ against a true $0.0868$ for the base engine and
$0.0668$ against $0.0660$ for the shipped model, the residual being within-bin forecast
variance):

\begin{table}[htbp]
\centering
\small
\begin{tabular}{lrrrr}
\toprule
model & reliability $\downarrow$ & resolution $\uparrow$ & uncertainty & $\approx$ Brier \\
\midrule
engine (raw conditional logistic) & 0.0231 & 0.0374 & 0.1022 & 0.0879 \\
shipped (top-5 warp + shrunk + tilt) & \textbf{0.0047} & 0.0401 & 0.1022 & 0.0668 \\
\bottomrule
\end{tabular}
\caption{Where the improvement came from. Reliability falls by a factor of five;
resolution barely moves.}
\label{tab:murphy}
\end{table}

This table is the most informative single result in the paper. The entire gain from the
tilt, the shrink and the room warp is a \textbf{calibration} gain, not a
\textbf{discrimination} gain. The model was always roughly as good at telling players
apart as it is now; it was simply saying the wrong numbers. That is worth knowing for
two reasons: it tells you the tilt is doing exactly what it was designed to do and
nothing sneaky besides, and it tells you where the remaining headroom is. If you want a
better engine, you need better \emph{resolution}, and no scalar correction will ever
supply it --- which is precisely the argument for the opponent-aware Monte Carlo model
of milestone 6.

\subsection{The backtest}
\label{sec:backtest}

\subsubsection*{Data}

One completed draft: Sleeper draft \code{1133489617308684288}, this user's real 2024
league draft. 12 teams, 15 rounds, 180 picks.

Era-correct prices: FFC's archive serves past years, and
\code{/api/v1/adp/half-ppr?teams=12\&year=2024} returns the 2024 market --- $178$
players over $906$ drafts, window 2024-08-31 to 2024-09-01. The rows carry the full
current schema (\code{adp}, \code{stdev}, \code{times\_drafted}, \code{high},
\code{low}), so per-player sigma, the shrink and the support floor are all directly
backtestable rather than only $\sigma_{\text{def}}$.

The draft ran on 2024-08-31, \emph{inside} that window. The tool checks and reports this
rather than assuming it.

The two 2025 drafts are \textbf{skipped, not scored}. FFC's archive has no
\code{year=2025} --- it returns ``No ADP data found''. Falling back to another season's
prices would silently invalidate every number, so the tool prints a note and drops them.

\subsubsection*{Vantages}

\begin{definition}[Vantage]
A \emph{vantage} is a triple (draft, seat, round): standing at seat $s$'s pick in round
$r$, looking forward to seat $s$'s pick in round $r+1$.
\end{definition}

For every seat $s \in \{1,\dots,12\}$ and every round $r \in \{1,\dots,14\}$ we take
$\text{from} = \mathrm{myPick}(r)$ and $\text{to} = \mathrm{myPick}(r+1)$ using the
engine's \emph{own} snake arithmetic with $\text{mySlot} = s$. That is $12 \times 14 =
168$ vantages.

\subsubsection*{The prediction set, the prediction, and the label}

At vantage $(\text{from}, \text{to})$, for each player $j$ on the era board:

\begin{itemize}
\item \textbf{Availability predicate.} Include $j$ if his actual draft position
$d_j \ge \text{from}$, or he went undrafted entirely. (If he was already gone, there is
nothing to predict.)
\item \textbf{Prediction.} $\hat p_j = $ the engine's \code{PSurviveAt} with
$\code{PickNo} = \text{from}$ and horizon $\text{to}$, using 2024 ADP and 2024 sigma.
This is literally the shipped function, called with a state whose \code{PickNo} is set;
it is not a re-implementation.
\item \textbf{Label.} $y_j = 1$ if $d_j \ge \text{to}$ or he went undrafted --- he really
was still sitting there.
\end{itemize}

That yields $\mathbf{15{,}923}$ labelled predictions.

\subsubsection*{The honest caveat: 12 seats multiply vantages, not evidence}

This must be said plainly because the prediction count invites the wrong conclusion.
Twelve seats over the same 180 picks gives twelve times the labelled data \emph{for
free}, and it is legitimate --- snake arithmetic is seat-agnostic, so each seat is a
genuinely different schedule of vantages over the same reality, and using all of them
removes any dependence on which seat we happened to pick.

But it is \textbf{not} twelve times the \emph{evidence}. Every seat is scoring the same
$180$ real picks. If Player X unexpectedly lasted to pick 90 in that draft, all twelve
seats record that same surprise, at overlapping horizons. The predictions are heavily
correlated and the \emph{effective} sample size is far smaller than $15{,}923$ ---
closer, in spirit, to one draft's worth of independent surprises.

Concretely, this means: do not compute a standard error from $n = 15{,}923$ and conclude
that a $-0.0009$ Brier difference is significant. It is not, and we do not claim it is.
The tool prints this caveat on every run.

\subsubsection*{The join must be exact-only}

Matching 2024 FFC names to Sleeper player IDs uses the project's normal matching
pipeline --- but with the fuzzy stage \emph{disabled}, and this is not fussiness.

The normal pipeline falls through to Levenshtein-distance matching within a position. A
2024 player who is missing from the \emph{current} active pool (retired, cut, out of the
league) would then silently map to a similarly-named active player --- whose ID never
appears in the 2024 pick list. He would therefore be labelled ``survived'' at every
single vantage. That is not a gap in the data, it is a \emph{bias}, and it points the
same direction every time: it inflates the observed survival rate and flatters every
model that predicts high.

So calibrate uses \code{Index.LookupExact} and prints anything it drops. On the 2024
board: $177$ of $178$ matched, $1$ dropped.

\subsubsection*{Leave-one-out for the room warp}

The room curve is built from this league's drafts --- and one of them is the draft being
scored. A warp built from the draft it is graded on has memorised the answer and would
report a skill it does not have.

So the gate builds the curve from the \emph{other two} drafts only (both 2025), and
scores it on 2024. That is leave-one-out cross-validation with $k = 1$ fold, which is
the only fold this data can cut, and \S\ref{sec:roomfail} and \S\ref{sec:threats} both
say so.

One further subtlety the implementation gets right and which is easy to miss: the warp
ranks players by ADP \emph{within a position}, so the rank input must cover \emph{every}
row on the era board, including the one name the exact-only join dropped. Otherwise a
missing name shifts every deeper player at that position up by one rank and warps him
with the wrong $\adp_{\text{room}}(P,k)$. Measured: the one dropped 2024 name
(a WR ranked 37th) would have mispriced the $28$ receivers behind him by $1.6$ picks on
average against a mean warp of $4.4$ --- an uncontrolled error inside the measurement the
whole phase rests on.

\subsubsection*{Solving the tilt at a vantage}

The tilt is not a per-player transform; it is one exponent solved over a whole vantage's
board at once. So a flat list of predictions is not enough to compute it, and each row
carries its vantage id. Crucially, when the results are broken out by segment (horizon,
position, ADP depth), the exponent for the ``13+ picks'' rows is the one solved over the
\emph{whole board} at that vantage --- re-solving over the filtered subset would invent a
draft in which only long-horizon players exist.

The backtest contains a second implementation of the tilt solver, because the engine's
own is unexported and reads its horizon off \code{NextPick()} --- and a backtest vantage
\emph{stands on} my own pick, so the engine's own solve degenerates to $c = 1$ and cannot
be driven to the right question from outside the package. Two copies of one model is a
liability, so it is checked rather than assumed: on every run, both solvers run over one
shared 60-player board and the maximum disagreement is printed (currently
$0.0\times10^{0}$), and a test fails the build if they ever diverge past $10^{-12}$.

\subsection{The baselines}

A model must beat all of these or the mathematics is theatre.

\begin{enumerate}
\item[(a)] \textbf{Constant.} Predict the global observed survival rate for everyone.
This is the ``no model at all'' floor and, by Theorem~\ref{thm:proper}, it scores exactly
the uncertainty term $\bar o(1-\bar o)$.
\item[(b)] \textbf{Flat sigma.} Replace every per-player $\sigma_j$ with
$\sigma_{\text{def}} = 6.0$. If this ties the engine, the whole stdev plumbing of
\S\ref{sec:sigma} is decoration.
\item[(c)] \textbf{Unconditional.} Drop the $S(q)/S(p)$ ratio and use the raw curve
$S_j(\text{to})$. If this ties the engine, conditioning on ``he is demonstrably still
here'' (\S\ref{sec:conditioning}) is decoration.
\end{enumerate}

%=======================================================================
\section{Results}
\label{sec:results}
%=======================================================================

All numbers below are the live output of \code{pick6 calibrate} on the 2024 draft:
$15{,}923$ predictions, $168$ vantages, observed survival rate $0.8844$.

\subsection{The model table}

\begin{table}[htbp]
\centering
\small
\begin{tabular}{lrrr}
\toprule
model & Brier $\downarrow$ & log-loss $\downarrow$ & mean predicted \\
\midrule
engine (per-player sigma)              & 0.0868 & 0.3007 & 0.8231 \\
engine + exactly-$N$ tilt              & 0.0677 & 0.2288 & 0.8734 \\
\midrule
4b: shrunk sigma                       & 0.0856 & 0.2895 & 0.8237 \\
4b: shrunk sigma + tilt                & 0.0670 & 0.2250 & 0.8734 \\
4b: support floor                      & 0.0868 & 0.3007 & 0.8231 \\
4b: support floor + tilt               & 0.0677 & 0.2288 & 0.8734 \\
4b: shrunk + floor + tilt              & 0.0670 & 0.2250 & 0.8734 \\
\midrule
3a: room-warped ADP (full depth)       & 0.0823 & 0.2867 & 0.8356 \\
3a: room warp + tilt                   & 0.0677 & 0.2364 & 0.8734 \\
3a: room warp + shrunk + tilt          & 0.0671 & 0.2327 & 0.8734 \\
\textbf{3a: top-5 warp + shrunk + tilt} & \textbf{0.0660} & \textbf{0.2222} & 0.8734 \\
\midrule
baseline: constant 0.884               & 0.1022 & 0.3580 & 0.8844 \\
baseline: sigma 6.0 flat               & 0.0825 & 0.2675 & 0.8330 \\
baseline: sigma 6.0 flat + tilt        & 0.0671 & 0.2233 & 0.8734 \\
baseline: unconditional                & 0.1111 & 0.5159 & 0.7934 \\
\bottomrule
\end{tabular}
\caption{Every model scored on the same $15{,}923$ predictions. Observed survival
$0.8844$. The bolded row is what the tool runs.}
\label{tab:models}
\end{table}

Five things this table says:

\begin{enumerate}
\item \textbf{The tilt is the single biggest win in the project}, by a distance. It moves
Brier from $0.0868$ to $0.0677$ ($-22\%$) and log-loss from $0.3007$ to $0.2288$
($-24\%$), and it pulls mean predicted survival from $0.8231$ to $0.8734$ against an
observed $0.8844$. A one-parameter correction solved from a budget constraint, with no
fitting to the labels whatsoever.

\item \textbf{Conditioning is load-bearing.} The unconditional baseline is the worst row
in the table --- worse than the constant predictor on Brier ($0.1111$ vs $0.1022$) and
catastrophically worse on log-loss ($0.5159$ vs $0.3580$). \S\ref{sec:conditioning} is
not a refinement, it is the difference between a model and an anti-model.

\item \textbf{Per-player sigma does not, on its own, beat a flat sigma --- and this is the
most uncomfortable row in the table.} Baseline (b) exists to test whether the whole stdev
pipeline of \S\ref{sec:sigma} earns its keep. It does not clear the bar cleanly. Untilted,
flat $\sigma = 6$ beats per-player sigma on \emph{both} metrics ($0.0825 / 0.2675$ against
$0.0868 / 0.3007$). Tilted, flat still edges it ($0.0671 / 0.2233$ against
$0.0677 / 0.2288$). Only once shrinkage is applied does per-player sigma draw level
($0.0670 / 0.2250$ --- better on Brier by $0.0001$, worse on log-loss by $0.0017$), and
only with the room warp on top does the full stack beat flat+tilt on both
($0.0660 / 0.2222$).

The honest reading: \emph{the per-player spread is real information, but the raw
conversion \eqref{eq:sigmaclamp} over-trusts it badly enough to be a net loss, and most of
what shrinkage buys is undoing that damage rather than adding skill.} The Jets-defense
case of \S\ref{sec:shrink} is not an anecdote --- it is the whole effect, and there are
$53$ players like him on the board. On this evidence, running the tool with
$\sigma \equiv 6.0$ and the tilt would cost you almost nothing. Per-player sigma stays
because it is right in principle, because the shipped stack does win, and because a
constant scale cannot possibly be right for both a $\varsigma = 0.4$ receiver and a
$\varsigma = 39.2$ kicker --- but ``the backtest demanded it'' would be false, and one
draft is not enough to settle it either way.

\item \textbf{The support floor does literally nothing.} Identical to four decimals with
and without, tilted and untilted. \S\ref{sec:negative} explains why, and it is
structural.

\item \textbf{The full-depth room warp loses; the top-5 warp wins.} $0.0671 / 0.2327$
against $0.0660 / 0.2222$, with the shrunk+tilt base at $0.0670 / 0.2250$.
\end{enumerate}

\subsection{Calibration, before and after}

\begin{table}[htbp]
\centering
\small
\begin{tabular}{lrrrrrr}
\toprule
bin & engine & +tilt & +shrunk & shipped & observed & $n$ \\
\midrule
0.0--0.1 & 0.0274 & 0.1535 & 0.1670 & 0.1715 & 0.3413 & 1383 \\
0.1--0.2 & 0.1481 & 0.4130 & 0.4154 & 0.4183 & 0.6432 & 426 \\
0.2--0.3 & 0.2495 & 0.5304 & 0.5269 & 0.5283 & 0.6520 & 319 \\
0.3--0.4 & 0.3475 & 0.6163 & 0.6082 & 0.6075 & 0.7021 & 292 \\
0.4--0.5 & 0.4491 & 0.7028 & 0.6934 & 0.6913 & 0.7190 & 306 \\
0.5--0.6 & 0.5496 & 0.7667 & 0.7581 & 0.7555 & 0.7419 & 279 \\
0.6--0.7 & 0.6504 & 0.8284 & 0.8211 & 0.8171 & 0.7710 & 310 \\
0.7--0.8 & 0.7527 & 0.8855 & 0.8800 & 0.8769 & 0.7995 & 414 \\
0.8--0.9 & 0.8544 & 0.9347 & 0.9309 & 0.9277 & 0.8678 & 643 \\
0.9--1.0 & 0.9931 & 0.9970 & 0.9967 & 0.9965 & 0.9842 & 11551 \\
\bottomrule
\end{tabular}
\caption{Bins are set by the \emph{engine's} prediction so all four models share the
same rows and are directly comparable within a row. Compare each model's column against
the observed column.}
\label{tab:relcompare}
\end{table}

Read the low bins first, because that is where the original model was broken. In the
bottom bin the engine said $0.027$ against an observed $0.341$; the tilted models say
$0.15$--$0.17$. Still under-confident, but the error has shrunk from a factor of $12$ to
a factor of $2$.

Now read the high bins, because that is the cost. In the $0.6$--$0.9$ bins the tilted
models \emph{over}-predict where the engine under-predicted: at bin $0.7$--$0.8$ the
engine said $0.753$ against $0.800$ observed (too low), the tilted model says $0.886$
(too high). The tilt is a global correction and it cannot fix the shape of the curve, only
its level; pushing the badly-wrong low bins up necessarily pushes the nearly-right high
bins past the mark. The trade is strongly worth it --- Table~\ref{tab:murphy}'s
reliability term falls five-fold --- but it is a trade, and it is visible here.

Table~\ref{tab:relship} is the shipped model graded on \emph{its own} bins, which is the
honest way to read one model's calibration rather than compare four.

\begin{table}[htbp]
\centering
\small
\begin{tabular}{lrrr}
\toprule
bin & mean predicted & observed & $n$ \\
\midrule
0.0--0.1 & 0.0392 & 0.1361 & 485 \\
0.1--0.2 & 0.1525 & 0.3575 & 372 \\
0.2--0.3 & 0.2503 & 0.4636 & 343 \\
0.3--0.4 & 0.3496 & 0.5937 & 315 \\
0.4--0.5 & 0.4494 & 0.6172 & 337 \\
0.5--0.6 & 0.5502 & 0.6619 & 349 \\
0.6--0.7 & 0.6515 & 0.7179 & 390 \\
0.7--0.8 & 0.7527 & 0.7225 & 454 \\
0.8--0.9 & 0.8546 & 0.8064 & 661 \\
0.9--1.0 & 0.9930 & 0.9788 & 12217 \\
\bottomrule
\end{tabular}
\caption{The shipped model, binned by its own prediction. Still under-confident below
$0.7$ and slightly over-confident above it, but nowhere near the original's
one-directional collapse.}
\label{tab:relship}
\end{table}

\subsection{Segments}

The tails are where the tool earns its keep, so they get reported separately.

\begin{table}[htbp]
\centering
\small
\begin{tabular}{lrrrr}
\toprule
segment & $n$ & observed & Brier (tilted $\to$ +shrunk) & log-loss \\
\midrule
\multicolumn{5}{l}{\emph{by horizon}}\\
$\le 6$ picks & 3804 & 0.9706 & $0.0255 \to 0.0254$ & $0.1044 \to 0.1041$ \\
7--12 picks   & 3919 & 0.9132 & $0.0629 \to 0.0625$ & $0.2176 \to 0.2161$ \\
13+ picks     & 8200 & 0.8307 & $0.0896 \to 0.0884$ & $0.2918 \to 0.2853$ \\
\midrule
\multicolumn{5}{l}{\emph{by position}}\\
QB  & 2447 & 0.9166 & $0.0962 \to 0.0970$ & $0.2961 \to 0.2966$ \\
RB  & 4181 & 0.8627 & $0.0586 \to 0.0589$ & $0.1863 \to 0.1861$ \\
WR  & 4752 & 0.8554 & $0.0675 \to 0.0658$ & $0.2494 \to 0.2450$ \\
TE  & 1742 & 0.8961 & $0.0884 \to 0.0849$ & $0.3281 \to 0.3101$ \\
K   & 1458 & 0.9328 & $0.0485 \to 0.0487$ & $0.1571 \to 0.1544$ \\
DEF & 1343 & 0.9285 & $0.0387 \to 0.0379$ & $0.1137 \to 0.1110$ \\
\midrule
\multicolumn{5}{l}{\emph{by ADP depth}}\\
rounds 1--3 (ADP $\le$ 36) & 789 & 0.5057 & $0.1077 \to 0.1073$ & $0.3390 \to 0.3357$ \\
rounds 4--8 (ADP $\le$ 96) & 4521 & 0.8354 & $0.0848 \to 0.0852$ & $0.2900 \to 0.2886$ \\
rounds 9+                  & 10613 & 0.9335 & $0.0574 \to 0.0562$ & $0.1945 \to 0.1896$ \\
\bottomrule
\end{tabular}
\caption{Segments, showing the effect of adding shrinkage to the tilted model.}
\label{tab:segments}
\end{table}

Observations:

\begin{itemize}
\item \textbf{Short horizons are easy and long ones are hard}, as they must be: at
$\le 6$ picks the observed survival rate is $0.97$ and the Brier is $0.0255$; at $13+$
picks the rate is $0.83$ and the Brier is $0.0884$. The tool is doing most of its real
work in the third row.

\item \textbf{Rounds 1--3 are the genuinely uncertain band}, and they show why a Brier
score must always be read against its own segment's base rate. Observed survival there is
$0.5057$ --- an almost perfect coin flip --- so that segment's uncertainty floor is
$0.5057 \times 0.4943 = 0.2500$, the full coin-flip ceiling. The model scores $0.1073$
against it: less than half. That is by some distance the strongest performance in the
table, even though $0.1073$ is numerically the \emph{worst} Brier of the three depth
rows. Judged against the global floor of $0.1022$ it looks like the model's weak spot; it
is the opposite. This is the segment where a draft is won, and the model is extracting
real signal from it.

\item \textbf{Quarterbacks are the worst-calibrated position} (predicted $0.841$ vs
observed $0.917$ under shrunk+tilt) and the top-5 warp improves them most
($0.0970 \to 0.0952$ Brier). \S\ref{sec:roomfail} explains why: it is the rank-versus-
player distinction.

\item \textbf{Tight ends benefit most from shrinkage} (Brier $0.0884 \to 0.0849$,
log-loss $0.3281 \to 0.3101$), which is what you would expect --- it is a shallow,
thinly-sampled position where raw spreads are noisiest.
\end{itemize}
%=======================================================================
\section{Negative results}
\label{sec:negative}
%=======================================================================

Three ideas were built, measured, and did not survive. They get their own section
because they are the best evidence that the measurement apparatus of
\S\ref{sec:calibration} is doing real work --- a backtest that only ever confirms is not
a backtest --- and because in two of the three cases the reason for failure is more
instructive than any of the successes.

Each subsection ends with what evidence would overturn it. That is not a formality: two
of the three are still in the codebase, exported and tested, one call site away from
being switched back on.

\subsection{The full-depth room warp}

Derived in full in \S\ref{sec:roomfail}; listed here so the three negatives sit together.

\textbf{What it was.} Blend every player's national ADP toward this room's own
rank-to-pick curve, at every depth. \textbf{Why it failed.} A ranked list runs deeper
than any finite draft, so past the room's appetite the rank-to-room-pick map has no data
--- and there are more players in that unreliable tail than at the reliable head. The
warp's mean \emph{blended} shift on quarterbacks flipped sign: $-1.7$ to $-3.7$ picks on
QB1 through QB5, but \textbf{$+11.0$ picks overall}, the opposite of the signal it was
built to capture. \textbf{The score.} Brier $0.0670 \to 0.0671$, log-loss $0.2250 \to 0.2327$;
both worse. Restricted to the top five at each position it wins instead
($0.0660 / 0.2222$), and that is what ships.

\textbf{Status.} \code{RoomCurve.EffectiveADP} is still exported, still tested, and still
scored as a row in the model table on every run --- deliberately, so the next reader
cannot assume the warp was never tried at depth. Its only caller is the backtest.

\textbf{What would overturn it.} A second scorable season, so the gate is not one fold;
or a curve built from enough drafts that the deep $k$ values have real support. This is a
small-sample-at-depth failure, not a wrong-idea failure, and ten seasons of this league's
drafts might well make the full-depth warp the better model.

\subsection{The support floor}
\label{sec:floorfail}

\textbf{What it was.} FFC reports \code{high}, the earliest pick anyone in its sample ever
took a player at. If my horizon is at or before \code{high}, then surviving to it means
surviving a window in which \emph{no observed draft ever removed him}. That is a zero-
event observation, and there is a clean bound for those.

\begin{definition}[Rule of three]
If an event occurred zero times in $n$ independent trials, an approximate $95\%$ upper
bound on its underlying rate is $3/n$ \cite{hanley1983}.
\end{definition}

\begin{proposition}[Where the 3 comes from]
If the true rate is $\pi$, the probability of observing zero events in $n$ trials is
$(1-\pi)^n$. Setting this to $0.05$ and solving,
\[
n\log(1-\pi) = \log 0.05 \quad\Longrightarrow\quad
\pi = 1 - e^{\log(0.05)/n} \approx \frac{-\log 0.05}{n} = \frac{2.996}{n}
\]
for small $\pi$, using $\log(1-\pi) \approx -\pi$. Hence $3/n$.
\end{proposition}

So the floor is
\begin{equation}
p \;\longleftarrow\; \max\!\left(p,\; 1 - \frac{3}{\max(n, 3)}\right)
\qquad\text{whenever } \text{horizon} \le \code{high},
\label{eq:floor}
\end{equation}
where $n = \code{times\_drafted}$. The $\max(n,3)$ keeps a two-draft player from
producing a negative floor: at $n=1$ the bound $1 - 3/1 = -2$ is the rule of three saying
it has nothing to offer, not a probability.

It only ever \emph{raises} a survival probability, never lowers one --- an observed early
pick is evidence a player \emph{can} go early, not evidence about how often.

\textbf{Why it failed, and the reason is beautiful.} It qualifies on $9{,}817$ of the
$15{,}923$ backtest rows --- $62\%$ of them. It \textbf{binds on zero}. Against the
shrunk curve it binds on two. Brier and log-loss are identical to four decimal places
with it and without it, tilted or untilted.

The reason is structural, not seasonal:

\begin{proposition}[The floor's window always sits where the curve is already near 1]
\label{prop:floordead}
\code{high} is a \emph{minimum} over drafts and \code{adp} is a \emph{mean} over the same
drafts. Hence $\code{high} \le \code{adp}$ always. Therefore the floor's applicability
window $\{\text{horizon} \le \code{high}\}$ lies entirely at or before the player's own
price --- and by \eqref{eq:S}, $S_j(p) \ge 1/2$ for all $p \le a_j$, with $S_j$ climbing
rapidly toward $1$ as $p$ falls below $a_j$.
\end{proposition}

\begin{proof}
The minimum of a finite set of positive numbers is at most their mean, so
$\code{high} \le \code{adp}$. Then $\text{horizon} \le \code{high} \le a_j$ gives
$(\text{horizon} - a_j)/\sigma_j \le 0$, so by \eqref{eq:S} the survival is at least
$1/(1+e^0) = 1/2$; the conditional version \eqref{eq:psurvive} is at least as large again,
since dividing by $S_j(p_{\text{now}}) \le 1$ only increases it.
\end{proof}

So the bound is competing against a curve that already reads near $1$ in the entire
region where the bound applies. It is dominated almost everywhere \emph{by construction}.
It was a good idea aimed at a real failure --- a tight curve inventing risk the sample
explicitly denies --- and the failure simply does not occur in the region where this
particular bound can speak.

Measured a second way, on the shipped 2026 board rather than the 2024 backtest: replayed
over all twelve seats and every round, it bites on $16$ of $34{,}104$ predictions, the
largest single move being $74.5\% \to 82.4\%$.

\textbf{Status.} Built, exported as \code{engine.SupportFloor}, unit-tested, graded on
every calibrate run, and \textbf{deliberately not wired in}. It is one call site away:
wrap \code{PSurviveAt}'s return. The ordering would matter --- the floor must come
\emph{before} the tilt, because the tilt balances a budget over the finished per-player
numbers, and flooring afterwards would hand back mass the tilt had just balanced and
break the exactly-$N$ property the tilt exists for.

\textbf{What would overturn it.} A source whose \code{high} is not a per-draft minimum
--- for instance a reported $5$th percentile of draft position, which can and does sit
\emph{after} the mean for a right-skewed player. Then Proposition~\ref{prop:floordead}
no longer applies and the window can land where the curve is genuinely uncertain.
Alternatively, if calibrate ever starts reporting a bind count worth grading, rewire it.

\subsection{The SigmaMin tuning artifact}
\label{sec:tuningartifact}

This one is a methodology lesson rather than a modelling one, and it is the most
transferable thing in this section.

\textbf{What happened.} \code{pick6 calibrate -tune} grid-searches
$\sigma_{\text{def}} \times \sigma_{\min} \times \sigma_{\max} \times n_0$ --- four axes,
$594$ combinations --- and \emph{today} it ranks them by the \textbf{tilted} Brier, which
is the model that ships. It did not always: before the tilt and the shrink existed there
was neither a tilted column nor an $n_0$ axis, so it was three axes ranked by the plain
Brier. In that state it recommended widening $\sigma_{\min}$ from $0.5$ to about $14$ ---
an enormous change, scoring $0.0690$ against the then-current $0.0868$.

Run \emph{after} the tilt, $\sigma_{\min} = 14$ scores $0.0690$ tilted as well --- which
is now \textbf{worse} than the shipped $0.0670$. (It is still the untilted optimum. Both
numbers reproduce on every run.)

\textbf{Why.} A very large $\sigma_{\min}$ flattens every tight player's curve, which
raises his survival probability. That is the same correction the tilt makes, arrived at
by a blunter route: widening the curve was compensating for the under-prediction bias of
Table~\ref{tab:premiscal}. Once the tilt removes that bias properly --- globally, with
one interpretable parameter, and without destroying the per-player spread information
that $\sigma_{\min}$ was flattening away --- the crutch becomes dead weight.

\textbf{The lesson.} \emph{Tune the model you ship, not a model you no longer run.} An
optimum found under one model configuration is not an optimum under another, and the two
can disagree not because one search was wrong but because a later change absorbed the
same error the earlier optimum was compensating for. Constants tuned against a
superseded pipeline are not neutral leftovers --- they are corrections for a bias that
has already been corrected, applied twice.

\subsubsection*{What the tilted grid says now, which is not nothing}

The tuner having been rebuilt to rank by the shipped metric, it is fair to ask what it
recommends. The answer is honest and slightly awkward, so here it is in full: the tilted
grid's optimum is \emph{not} the shipped combination. Ranked by tilted Brier, the top of
the table reads

\begin{center}
\small
\begin{tabular}{rrrrrr}
\toprule
$\sigma_{\text{def}}$ & $\sigma_{\min}$ & $\sigma_{\max}$ & $n_0$ & untilted Brier & tilted Brier \\
\midrule
4.5 & 8.0 & 10.0 & 200 & 0.0750 & \textbf{0.0667} \\
4.5 & 4.0 & 15.0 & 50  & 0.0833 & \textbf{0.0667} \\
4.5 & 8.0 & 15.0 & 50  & 0.0749 & 0.0668 \\
\midrule
6.0 & 0.5 & 25.0 & 25  & 0.0856 & 0.0670 \quad (shipped) \\
\bottomrule
\end{tabular}
\end{center}

\noindent
So the overturn condition stated below \emph{is} partially met: with the tilt in place the
grid finds $0.0667$ against the shipped $0.0670$. Three things about that margin, all
checkable in the tool's own output.

\begin{enumerate}
\item \textbf{It is a flat-sigma model in disguise.} Blend at $n_0 = 200$ and the prior
swamps every player's own measurement (median \code{times\_drafted} on the 2024 board is
$54$); clamp the result to $[8, 10]$ and there is almost nothing left to vary. Measured:
that row gives $\sigma$ running $8.00$ / $8.00$ / $8.99$ (min / median / max) with
$154$ of $178$ players pinned exactly at the floor. The second row is milder and still
pins $81$ of $178$. The grid is not discovering a better per-player scale --- it is
rediscovering flat sigma, which the results table already grades directly: ``baseline:
sigma 6.0 flat + tilt'' scores $0.0671$. A $0.0667$ and a $0.0671$ that are the same
model in different clothes are not evidence for moving four constants.

\item \textbf{It sits on a grid edge.} $n_0 = 200$ is the top of the shrink axis, and
\code{edgeWarning} prints ``the winner sits on the edge of the shrink axis --- widen the
grid before believing it'' on every run. An optimum at the boundary is the search running
out of road, not a minimum.

\item \textbf{The margin is $0.0003$ on one fold.} \S\ref{sec:threats} says why that is
not a number to move constants on. The room warp ships on $0.0009$, which is three times
larger, and it additionally cleared a cutoff sweep showing a plateau and a per-position
gate showing no regressions. This row has neither.
\end{enumerate}

\textbf{Outcome.} No constant was moved. $\sigma_{\min}$ is still $0.5$, and the shipped
row is printed in the grid on every run precisely so this decision stays visible rather
than becoming folklore.

\textbf{What would overturn it.} A second scorable season on which the same region of the
grid still wins, or an interior winner that is not simply flat sigma wearing four
constants. The tuner still exists and still prints; it never auto-writes. Printing rather
than writing is the whole reason any of this was catchable.

%=======================================================================
\section{Threats to validity}
\label{sec:threats}
%=======================================================================

A number without its caveat is a lie. These are ordered roughly by how much they should
worry you.

\paragraph{1. One draft.}
The entire backtest is a single completed draft. There is exactly one season for which
this league's draft and the corresponding era ADP both exist. Every claim in
\S\ref{sec:results} is therefore a claim about one night in August 2024. The room warp's
margin in particular --- $-0.0009$ Brier, $-0.0028$ log-loss --- is small enough that a
single unusual draft could produce it by chance, and we cannot rule that out with the
data we have.

\paragraph{2. Twelve seats are correlated.}
$15{,}923$ predictions sounds like a lot. It is $168$ vantages over $180$ real picks, and
all twelve seats are scoring the same picks. The effective sample is far smaller than the
count suggests. Do not compute a standard error from $n = 15{,}923$.

\paragraph{3. The ADP snapshot postdates the draft --- and this is the optimistic case.}
FFC serves the trailing snapshot of a season's draft week. This draft ran on 2024-08-31,
inside the 2024-08-31\,--\,2024-09-01 window, so the prices we score against are the
prices the room actually had. But those prices already know what the room knew: late
injuries, camp news, the whole final week of information. Read every number in
\S\ref{sec:results} as \textbf{the model running with final prices} --- the best case,
not a typical case. On draft day the tool runs on a snapshot that is at most a day old,
which is close but not identical; the tool prints its own staleness in the footer for
exactly this reason.

\paragraph{4. 2025 is missing from the archive.}
Two of this user's three completed drafts contribute nothing to survival scoring, because
FFC's archive has no \code{year=2025}. They do contribute to the room curve, which reads
only pick order and is therefore era-immune. Recheck the archive near draft day --- it is
one request --- and if 2025 has landed, re-gate \S\ref{sec:roomfail} immediately. That is
the single highest-value follow-up in the project.

\paragraph{5. The room curve is three drafts.}
$552$ picks across three nights of one league, two of them from 2025 and one from 2024.
The pseudo-count of \eqref{eq:warp} is an explicit acknowledgement of this ($w = 0.6$ at
$n = 3$), and the top-5 restriction is a second one. Neither makes three drafts into a
sample. And the room's membership changes between seasons, which nothing here models at
all.

\paragraph{6. Values are imported; the project makes no projections.}
Every $v_j$ comes from FantasyCalc, every $a_j$ and $\varsigma_j$ from FFC, every tier
from a human or from a value gap. If the imported value curve is wrong about a player,
the engine is wrong about him in exactly the same way and has no independent means of
noticing. This is a deliberate scope decision --- ``no projections of our own'' is the
project's first non-goal --- but it means the engine's ceiling is the value source's
ceiling. Urgency can only tell you what waiting costs \emph{in the value source's own
units}.

\paragraph{7. Independence is assumed twice and is false both times.}
Theorem~\ref{thm:ebest} and \eqref{eq:tierhold} both assume players survive
independently. They do not: a run on running backs removes several at once, and the
manager who takes one is thereby not taking another. The tilt corrects the aggregate
consequence of this (the removal budget) but not the correlation structure. Modelling
that properly is exactly what the Monte Carlo opponent model of milestone 6 is for, and
until it lands the engine's tier-hold probabilities in particular should be read as
somewhat too confident in both directions.

\paragraph{8. Need weights are unmeasured.}
$1.0 / 0.6 / 0.25$ and the endgame slack of $0.5$ are judgement, not measurement. Nothing
in the backtest touches them, because the backtest grades survival and need does not enter
survival. They could be materially wrong without any of the numbers in this paper moving.

\paragraph{9. Per-player sigma is carried on principle, not on evidence.}
See point 3 of \S\ref{sec:results}: on this backtest a flat $\sigma = 6.0$ with the tilt
scores $0.0671 / 0.2233$ against the shrunk per-player stack's $0.0670 / 0.2250$ --- a
tie at best, and untilted the flat model wins outright. The shipped model beats it only
once the room warp is added, and that rests on the same single fold. \S\ref{sec:tuningartifact}
points the same way from a second direction: the tilted tuning grid's optimum is a
configuration that pins $154$ of $178$ players to one $\sigma$. If a second scorable
season arrives and per-player sigma still fails to beat a constant, the honest response is
to delete a lot of \S\ref{sec:sigma} rather than defend it. What keeps it in the meantime
is that a single scale cannot be right for both a $\varsigma = 0.4$ receiver and a
$\varsigma = 39.2$ kicker, and being wrong about those two is the whole job.

\paragraph{10. The two-pick plan's first leg is knowingly optimistic.}
\S\ref{sec:plan} prices the first leg at today's best available even when the pick is
seventeen away. The comparison between pairs survives it; the absolute score does not,
which is why the score is not displayed.

\paragraph{11. \code{TierMaxSize} does not reach derived tiers.}
\S\ref{sec:tierhold} states the scope and this is the consequence: \code{subdivide} runs
only on blocks a rankings file drew, so on a rankings-free board nothing caps a
value-gap tier's size. Today that leaves an $11$-man RB tier 2 near the top of the board
and a $29$-man WR tier 9 at the bottom, and \eqref{eq:tierhold} over eleven or twenty-nine
players is never going to fall below \code{TierHoldWarn}. The cost is silence, not a false
alarm --- those groups simply never raise a cliff --- and it is invisible in every number
in \S\ref{sec:results}, because the backtest grades survival and tiers do not enter
survival. It is a real gap all the same, and the fix is one call to \code{subdivide} on
\code{deriveBlocks}' output.

%=======================================================================
\section{Constants}
\label{sec:constants}
%=======================================================================

Every tuneable number, its value, where it lives, and the argument behind it. The
engine's constants are in \code{internal/engine/tuning.go}; the fetch-time duplicates are
in \code{internal/adp}, and that duplication is deliberate --- the \code{adp} package
precomputes sigma and tiers without importing the engine.

\begingroup
\footnotesize
\setlength{\LTleft}{0pt}\setlength{\LTright}{0pt}
\setlength{\LTpre}{4pt}\setlength{\LTpost}{4pt}
\renewcommand{\arraystretch}{1.05}
\raggedright
\begin{longtable}{@{}p{0.24\textwidth}p{0.07\textwidth}p{0.595\textwidth}@{}}
\toprule
\textbf{constant} & \textbf{value} & \textbf{where it comes from} \\
\midrule
\endfirsthead
\toprule
\textbf{constant} & \textbf{value} & \textbf{where it comes from} \\
\midrule
\endhead
\multicolumn{3}{@{}l}{\emph{Survival curve --- \code{engine/tuning.go}}}\\[0.2em]
\code{SigmaDefault} & 6.0 & Median $\varsigma$ across the 2026 half-PPR pool is $10.3$; $10.3/1.8138 = 5.68$, so $6.0$ is within six percent of the median player's own scale. Well calibrated there, wrong at both tails --- which is why it is only the fallback. On this board it never fires at all: FFC reports a spread for every row. \\
\code{SigmaFromStdev} & 1.8138 & $\pi/\sqrt3$, exactly. Proposition~\ref{prop:logisticvar}. Not tuneable; it is a theorem. \\
\code{SigmaMin} & 0.5 & Judgement. A per-pick hazard of $86\%$ is already ``he is gone''; below this the number stops meaning anything. Binds on $3$ of $203$ players. The \code{-tune} grid wanted $14$ here before the tilt and $8$ after it; \S\ref{sec:tuningartifact} is why it got neither. \\
\code{SigmaMax} & 25.0 & Judgement. $\varsigma \ge 45.3$ maps here; a player the market disagrees about by more than three rounds is already ``no idea''. Binds on \emph{nobody} on the 2026 board, whose widest spread is $\varsigma = 39.2$. \\
\code{SurviveThreshold} & 0.5 & The ``safe to wait'' cutoff --- \emph{display only} since urgency became continuous. No longer any part of the maths. \\
\code{SurvivalExpClamp} & 30.0 & $\log(1+e^{30}) - 30 = 9.4\times10^{-14}$, so the linear branch of softplus is exact to $10^{-14}$. See \S\ref{sec:logspace} for why this must never be applied to $S$ itself. \\
\code{UndraftedADP} & 999.0 & Off the drafted radar. Live feeds register handcuffs and rookies no ADP source ranked. \\
\code{FallerSigmas} & 1.0 & Judgement. One of the player's \emph{own} scale units past his price. Measuring in his own $\sigma$ is the load-bearing part; the $1.0$ is taste. \\
\code{RuleOfThree} & 3.0 & $-\log(0.05) = 2.996$. \S\ref{sec:floorfail}. Currently unreachable --- \code{SupportFloor} is not wired in. \\
\addlinespace
\multicolumn{3}{@{}l}{\emph{The tilt --- \code{engine/tuning.go}}}\\[0.2em]
\code{TiltCMax} & 64.0 & Bracket $[1/64, 64]$ for the bisection of Theorem~\ref{thm:tilt}. On real boards $c$ solves to $0.29$--$1.12$ and neither clamp ever fires: $0$ of $168$ backtest vantages, $0$ of $1{,}969$ mock frames. Wide enough to be irrelevant, which is the correct state for a clamp. \\
\code{TiltTol} & $10^{-6}$ & Bracket width at termination; $\approx 26$ halvings. \\
\addlinespace
\multicolumn{3}{@{}l}{\emph{Expected best / tiers --- \code{engine/tuning.go}}}\\[0.2em]
\code{EBestEpsilon} & $10^{-6}$ & Truncation of \eqref{eq:ebest}; omits at most $\varepsilon v_1 \approx 0.01$ value points on the shipped board. Proved in \S\ref{sec:ebest}. \\
\code{TierHoldWarn} & 0.5 & Judgement. Amber when the tier is a coin flip to survive. \\
\code{TierHoldCliff} & 0.15 & Judgement. Red when it almost certainly will not. \\
\code{TierDropPct} & 0.10 & Relative value drop that breaks a derived tier. Produces recognisable tiers on the real board --- with no rankings file, 7 QB, 6 RB, 9 WR and 4 TE tiers. \\
\code{TierFloorPct} & 0.015 & Absolute floor on that drop, as a fraction of the top value \emph{of the block being derived} --- the position's best when no rankings file is in play, the best uncovered player otherwise. Stops the flat tail shattering into singletons. \\
\code{TierMaxSize} & 8 & A published graphic put 21 receivers in one tier; a tier that large can never fire a cliff, so the board would never warn about receivers at all. Applies to rankings-file blocks only --- see \S\ref{sec:threats} item 11. \\
\code{TierMinSize} & 2 & Never manufacture a one-player tier --- that is a permanent ``last one'' alarm, the exact failure that got FantasyCalc's global tiers rejected. Singletons a human drew are left alone. \\
\addlinespace
\multicolumn{3}{@{}l}{\emph{Need --- \code{engine/tuning.go}}}\\[0.2em]
\code{NeedStarter} & 1.0 & Judgement, unmeasured. See \S\ref{sec:threats} item 8. \\
\code{NeedFlex} & 0.6 & Judgement, unmeasured. \\
\code{NeedBench} & 0.25 & Judgement, unmeasured. \\
\code{EndgameSlack} & 0.5 & Judgement. At $L = U+1$ a bench flier is still affordable but spends the only spare pick. At $L = U$ the multiplier is $0$, which is not tuneable --- it is the arithmetic of having no spare picks. \\
\code{KDefLastRounds} & 3 & \textbf{Policy, not a model.} What we are willing to recommend. \\
\code{OpponentKDefLastRounds} & 6 & \textbf{Measurement.} This room takes its first kicker in round 10 and first defense in round 11. Used by engine v2's opponent model only. Two constants because they are two questions. \\
\addlinespace
\multicolumn{3}{@{}l}{\emph{Runs, byes, freshness --- \code{engine/tuning.go}}}\\[0.2em]
\code{RunWindow} & 6 & Judgement. Half a round of a 12-team draft. \\
\code{RunThreshold} & 4 & Judgement. Four of six is unmistakably a run. \\
\code{ByeConflictThreshold} & 3 & Two starters sharing a bye is unremarkable across nine slots; three is a week you cannot field. \\
\code{StaleADPHours} & 24 & FFC recomputes daily, so a board over a day old has definitionally missed a refresh. \\
\code{ADPWindowStaleDays} & 2 & The source's own window, not our fetch timestamp: a cache-only fetch moves the timestamp and nothing else. \\
\code{NewsFreshHours} & 48 & Long enough to survive a draft-morning fetch about Friday-night news; short enough not to decorate half the board. \\
\code{TripwireSlack} & 2 & \code{low} is a maximum over the drafts FFC actually saw that player in --- $5$ to $215$ of them on the 2026 board, not the window's $1{,}187$. One pick behind it is a rounding difference; two means the room is genuinely behind every draft in his own sample. \\
\addlinespace
\multicolumn{3}{@{}l}{\emph{Value fallback --- \code{engine/tuning.go}}}\\[0.2em]
\code{ValueBase} & 250.0 & Only used when no source values anybody: $v = 250e^{-\mathrm{rank}/40}$. Never mixed with imported values --- \S\ref{sec:kdef} shows the 24-fold scale mismatch that would cause. \\
\code{ValueDecay} & 40.0 & As above. Chosen to be convex, not fitted. \\
\addlinespace
\multicolumn{3}{@{}l}{\emph{Fetch-time --- \code{internal/adp}}}\\[0.2em]
\code{ShrinkPseudoDrafts} & 25.0 & $n_0$ in \eqref{eq:shrinkvar}. The 2024 pool's median player was taken in $54$ drafts and the thinnest in $11$ (2026: $49$ and $5$); at $n_0 = 25$ the median player's own number outvotes the prior 2:1 and a thin player lands most of the way onto it. \\
\code{RoomWarpPseudo} & 2.0 & $w = n/(n+2)$, so three drafts buy $60\%$ of the warp and a depth nobody reached buys none. Deliberately timid. \\
\code{RoomWarpTopK} & 5 & \textbf{Measured}, and the one constant in this table chosen by the backtest. Cutoff sweep: $K=1$ $0.0665$, $K=3$ $0.0662$, $K=5$ $0.0660$, $K=8$ $0.0661$, $K=12$ $0.0674$, uncapped $0.0671$. A plateau at 4--5. \S\ref{sec:roomfail}. \\
\code{SigmaDefault}, \code{SigmaFromStdev}, \code{SigmaMin}, \code{SigmaMax}, \code{TierDropPct}, \code{TierFloorPct}, \code{TierMaxSize}, \code{TierMinSize} & --- & Deliberate duplicates of the engine values, so \code{fetch} can precompute sigma and tiers without importing the engine. Keep them in sync by hand; there are tests. \\
\bottomrule
\end{longtable}
\endgroup

The honest summary of this table: \textbf{exactly one constant was chosen by
measurement} (\code{RoomWarpTopK}), one is a theorem (\code{SigmaFromStdev}), several
are numerical-analysis facts with proofs attached (\code{SurvivalExpClamp},
\code{EBestEpsilon}, \code{TiltTol}, \code{RuleOfThree}), a handful are grounded in a
measured property of the data without being optimised against it
(\code{SigmaDefault}, \code{ShrinkPseudoDrafts}, \code{OpponentKDefLastRounds},
\code{TierMaxSize}), and the rest are judgement. That distribution is worth seeing
plainly. The parts of the engine that were graded are graded; the parts that were not are
labelled.
%=======================================================================
\section{From formula to code}
\label{sec:codemap}
%=======================================================================

This is the page to read with the repository open. Every numbered equation above, and
every named idea, mapped to the Go function that implements it. Function names are as
they appear in the source; lowercase names are unexported.

\begingroup
\footnotesize
\setlength{\LTleft}{0pt}\setlength{\LTright}{0pt}
\setlength{\LTpre}{4pt}\setlength{\LTpost}{4pt}
\renewcommand{\arraystretch}{1.05}
\raggedright
\begin{longtable}{@{}p{0.095\textwidth}p{0.205\textwidth}p{0.27\textwidth}p{0.35\textwidth}@{}}
\toprule
\textbf{eq.} & \textbf{what it is} & \textbf{function} & \textbf{file} \\
\midrule
\endfirsthead
\toprule
\textbf{eq.} & \textbf{what it is} & \textbf{function} & \textbf{file} \\
\midrule
\endhead

\multicolumn{4}{@{}l}{\emph{Snake arithmetic}}\\[0.2em]
\eqref{eq:round} & round of a pick & \code{State.Round} & \code{engine/state.go} \\
\eqref{eq:index} & index within round & \code{State.IndexInRound} & \code{engine/state.go} \\
\eqref{eq:slotat} & seat picking at $p$ & \code{State.SlotAt} & \code{engine/state.go} \\
\eqref{eq:mypick} & my pick in round $r$ & \code{State.MyPick}, \code{State.slotPosition} & \code{engine/state.go} \\
\eqref{eq:nextpick} & $q_1$ & \code{State.NextPick} & \code{engine/state.go} \\
\eqref{eq:following} & $q_2$ & \code{State.FollowingPick} & \code{engine/state.go} \\
\eqref{eq:picksuntil} & $N$ & \code{State.PicksUntilMine} & \code{engine/state.go} \\
--- & my next $n$ picks & \code{State.MyUpcomingPicks} & \code{engine/state.go} \\
--- & feed cross-check of \eqref{eq:slotat} & \code{State.ApplyRemote} & \code{engine/state.go} \\

\addlinespace
\multicolumn{4}{@{}l}{\emph{Survival}}\\[0.2em]
\eqref{eq:sigmafromstdev}, \eqref{eq:sigmaclamp} & $\varsigma \to \sigma$, clamped & \code{adp.Sigma} & \code{adp/aggregate.go} \\
\eqref{eq:S}, \eqref{eq:psurvive}, \eqref{eq:logspaceratio} & conditional survival & \code{State.PSurviveAt} & \code{engine/urgency.go} \\
--- & the same at horizon $q_1$ & \code{State.PSurvive} & \code{engine/urgency.go} \\
--- & $\log(1+e^x)$ with the clamp & \code{softplus} & \code{engine/urgency.go} \\
--- & which ADP the curve reads & \code{Player.price} & \code{engine/urgency.go} \\
\eqref{eq:falling} & falling flag & \code{State.Falling} & \code{engine/urgency.go} \\
\eqref{eq:prior24}, \eqref{eq:prior26} & OLS fit of $\varsigma$ on $a$ & \code{adp.FitStdevPrior}, \code{StdevPrior.At} & \code{adp/aggregate.go} \\
\eqref{eq:shrinkvar} & variance shrinkage & \code{adp.Blend}, \code{StdevPrior.Shrink} & \code{adp/aggregate.go} \\
--- & pool-wide shrink pass & \code{adp.ShrinkSigma} & \code{adp/aggregate.go} \\
\eqref{eq:floor} & rule-of-three floor (\textbf{not wired in}) & \code{engine.SupportFloor} & \code{engine/urgency.go} \\

\addlinespace
\multicolumn{4}{@{}l}{\emph{The tilt}}\\[0.2em]
\eqref{eq:oppbefore} & $N(\text{at})$ & \code{State.opponentPicksBefore} & \code{engine/expect.go} \\
\eqref{eq:tilt} & solve $f(c) = N$ by bisection & \code{State.survivalTilt} & \code{engine/expect.go} \\
--- & $\tilde p_j$ at horizon $q_1$ & \code{State.PSurviveTilted} & \code{engine/expect.go} \\
--- & backtest's second solver & \code{solveTilt}, \code{tilted} & \code{cmd/pick6/calibrate\_score.go} \\
--- & the two solvers agree & \code{tiltSelfCheck} & \code{cmd/pick6/calibrate\_score.go} \\

\addlinespace
\multicolumn{4}{@{}l}{\emph{Expected best survivor, urgency, need}}\\[0.2em]
\eqref{eq:ebest} & $\mathbb{E}[M_P]$ & \code{State.EBest}, \code{State.ebest} & \code{engine/expect.go} \\
--- & the sum itself, over slices & \code{expectedBest} & \code{engine/expect.go} \\
--- & modal best survivor (display) & \code{State.BestLater} & \code{engine/urgency.go} \\
\eqref{eq:urgency} & urgency & \code{State.Urgency} & \code{engine/urgency.go} \\
--- & ``safe to wait'' & \code{State.SafeToWait} & \code{engine/urgency.go} \\
\eqref{eq:need} & need weights & \code{State.Need}, \code{State.needFrom} & \code{engine/state.go} \\
--- & need after a hypothetical pick & \code{State.NeedAfter} & \code{engine/state.go} \\
Prop.~\ref{prop:greedy} & greedy slot filling & \code{State.FilledSlots}, \code{State.fillSlots} & \code{engine/state.go} \\
--- & K/DEF suppression & \code{State.Suppressed} & \code{engine/state.go} \\
\eqref{eq:endgame} & endgame slack & \code{endgameSlack}, \code{State.MustFillStarters} & \code{engine/state.go} \\
--- & board group ordering & \code{Board.bestAvailable} & \code{ui/board.go} \\

\addlinespace
\multicolumn{4}{@{}l}{\emph{Tiers, cliffs, runs}}\\[0.2em]
\eqref{eq:tierhold} & tier-hold probability & \code{State.TierHold} & \code{engine/detect.go} \\
--- & which horizon it prices to & \code{State.TierHoldPick} & \code{engine/detect.go} \\
--- & cliff level from the hold & \code{State.Cliff} & \code{engine/detect.go} \\
--- & remaining / total in a tier & \code{State.TierRemaining}, \code{State.TierSize} & \code{engine/state.go} \\
--- & run detection & \code{State.DetectRun} & \code{engine/detect.go} \\
--- & tier assignment & \code{adp.AssignTiers}, \code{assignPositionTiers} & \code{adp/aggregate.go} \\
--- & oversize-block subdivision & \code{subdivide}, \code{biggestDrop} & \code{adp/aggregate.go} \\
--- & derived tiers, runt merge & \code{deriveBlocks}, \code{mergeRuntTiers} & \code{adp/aggregate.go} \\

\addlinespace
\multicolumn{4}{@{}l}{\emph{Two-pick lookahead}}\\[0.2em]
\eqref{eq:plan} & the pair score and argmax & \code{State.BestPlan} & \code{engine/plan.go} \\
--- & rendering both legs by pick & \code{Board.planCopy}, \code{Board.planLine} & \code{ui/board.go} \\

\addlinespace
\multicolumn{4}{@{}l}{\emph{Room warp}}\\[0.2em]
\eqref{eq:roomcurve} & build $\adp_{\text{room}}$ & \code{adp.BuildRoomCurve} & \code{adp/room.go} \\
--- & held-out build for the gate & \code{adp.RoomCurveOf} & \code{adp/room.go} \\
--- & picks $\to$ positions in order & \code{adp.RoomDraftOf} & \code{adp/room.go} \\
--- & read one point of the curve & \code{RoomCurve.At}, \code{RoomCurve.Depth} & \code{adp/room.go} \\
\eqref{eq:warp} & the blend & \code{RoomCurve.Warp} & \code{adp/room.go} \\
--- & whole board, full depth (\textbf{loses the gate}) & \code{RoomCurve.EffectiveADP} & \code{adp/room.go} \\
--- & whole board, top $K$ (\textbf{ships}) & \code{RoomCurve.EffectiveADPTopK} & \code{adp/room.go} \\
--- & the three monotonicity guards & \code{RoomCurve.effectiveADP} & \code{adp/room.go} \\
--- & apply it at board load & \code{roomWarp}, \code{loadBoard} & \code{cmd/pick6/mock.go} \\
--- & read the cached drafts & \code{adp.CachedRoomDrafts}, \code{LoadRoomDrafts} & \code{adp/room.go} \\

\addlinespace
\multicolumn{4}{@{}l}{\emph{Value over replacement, K/DEF}}\\[0.2em]
\eqref{eq:vor} & $\mathrm{vor}$ & \code{State.VOR} & \code{engine/vor.go} \\
--- & $v_{\mathrm{rep}}(P)$ & \code{State.Replacement} & \code{engine/vor.go} \\
--- & $\mathrm{dem}(P)$, and its lineup fallback & \code{adp.PositionDemand}, \code{State.demandAt} & \code{adp/demand.go}, \code{engine/vor.go} \\
\S\ref{sec:kdef} & rank-anchored K/DEF values & \code{adp.AnchorKDefValues} & \code{adp/aggregate.go} \\

\addlinespace
\multicolumn{4}{@{}l}{\emph{Backtest}}\\[0.2em]
\S\ref{sec:backtest} & the whole command & \code{runCalibrate} & \code{cmd/pick6/calibrate.go} \\
--- & era board + exact-only join & \code{loadEraBoard}, \code{Index.LookupExact} & \code{cmd/pick6/calibrate.go}, \code{rankings/match.go} \\
--- & vantages, predictions, labels & \code{walk} & \code{cmd/pick6/calibrate.go} \\
Thm.~\ref{thm:proper} & Brier and log-loss & \code{scoreOf}, \code{clampP} & \code{cmd/pick6/calibrate\_score.go} \\
Thm.~\ref{thm:murphy} & reliability tables & \code{reliability} & \code{cmd/pick6/calibrate\_score.go} \\
--- & horizon / position / depth cuts & \code{segments}, \code{segmentTable}, \code{depthSegs} & \code{cmd/pick6/calibrate\_score.go} \\
--- & the baselines & \code{modelConstant}, \code{modelFlatSigma}, \code{modelUnconditional} & \code{cmd/pick6/calibrate\_score.go} \\
--- & 4b's two models & \code{modelShrunkSigma}, \code{modelSupportFloor} & \code{cmd/pick6/calibrate\_score.go} \\
\S\ref{sec:floorfail} & how often the floor bites & \code{floorReach} & \code{cmd/pick6/calibrate\_score.go} \\
\S\ref{sec:roomfail} & the 3a gate, per position & \code{roomGate}, \code{roomPosGate}, \code{shiftByPos} & \code{cmd/pick6/calibrate\_room.go} \\
\S\ref{sec:tuningartifact} & the sigma grid search & \code{tuneSigma}, \code{sigmaModel} & \code{cmd/pick6/calibrate\_tune.go} \\
--- & sigma sweeps with parameters & \code{psurvive}, \code{sigmaOf} & \code{cmd/pick6/calibrate\_score.go} \\
\bottomrule
\end{longtable}
\endgroup

One row in that table is a deliberate duplication rather than an oversight:
\code{solveTilt} in the backtest reimplements \code{State.survivalTilt}, for the reason
given at the end of \S\ref{sec:backtest}, and \code{tiltSelfCheck} is what keeps the two
honest.

%=======================================================================
\section{What is next}
%=======================================================================

Three things, in the order they should happen.

\paragraph{Re-gate when the archive grows.} FFC's \code{year=2025} still answers ``No ADP
data found''. The moment it does not, there is a second scorable season, the room warp's
one-fold gate becomes two folds, and \S\ref{sec:roomfail} gets an actual answer rather
than a defensible guess. This is one HTTP request and it is by far the highest-value
follow-up here.

\paragraph{Opponent-aware survival (milestone 6).} Everything in
\S\ref{sec:survival}--\S\ref{sec:tilt} treats the picks between mine as random draws
from ``drafters in general''. Sleeper hands us the actual rosters. A Monte Carlo that
simulates the specific eleven managers --- their unfilled slots, their measured
tendencies, their autodraft rates --- replaces the survival function and nothing else:
$\mathbb{E}[M_P]$, urgency, tier-hold and the banners all consume it unchanged.
Table~\ref{tab:murphy} says exactly why this is the right next move: the remaining
headroom is in \emph{resolution}, and no scalar correction can supply resolution.

\paragraph{Multi-pick lookahead (milestone 7).} \S\ref{sec:plan}'s honest simplification
--- pricing the first leg at today's best available --- goes away once a rollout is
simulating the draft forward anyway. Search over \emph{my} picks inside it, not just
over opponents'.

Two things that are deliberately \emph{not} next: an ADP over/under-performance model
(the blocker is labels --- it needs several seasons of historical ADP joined to actual
finish, and no source we have serves history), and any attempt to make projections of our
own. The first is parked on a data problem. The second is a scope decision and it is not
being revisited.

%=======================================================================
\begin{thebibliography}{9}
\addcontentsline{toc}{section}{References}

\bibitem{brier1950}
G.~W. Brier.
\emph{Verification of forecasts expressed in terms of probability.}
Monthly Weather Review, 78(1):1--3, 1950.
The original proposal of the mean squared error of a probability forecast as a
verification score. Brier's own formulation was for multi-category forecasts and
differs from the binary version used here by a factor of two; the binary
$\frac{1}{N}\sum(q_i - y_i)^2$ is what everyone now means by ``Brier score''.

\bibitem{murphy1973}
A.~H. Murphy.
\emph{A new vector partition of the probability score.}
Journal of Applied Meteorology, 12(4):595--600, 1973.
The three-way decomposition of the Brier score into reliability, resolution and
uncertainty, used here as Theorem~\ref{thm:murphy}. The decomposition is exact when
forecasts take finitely many distinct values and bins group identical forecasts;
binning continuous forecasts, as we do, leaves a within-bin variance residual.

\bibitem{hanley1983}
J.~A. Hanley and A.~Lippman-Hanley.
\emph{If nothing goes wrong, is everything all right? Interpreting zero numerators.}
Journal of the American Medical Association, 249(13):1743--1745, 1983.
The ``rule of three'': with zero events observed in $n$ trials, an approximate upper
$95\%$ confidence bound on the underlying rate is $3/n$. The short derivation is
reproduced in \S\ref{sec:floorfail}. The result is older than this paper and is often
attributed more diffusely; this is the reference that made it standard practice.

\bibitem{efronmorris1975}
B.~Efron and C.~Morris.
\emph{Data analysis using Stein's estimator and its generalizations.}
Journal of the American Statistical Association, 70(350):311--319, 1975.
The practical case for shrinkage: many noisy estimates pulled toward a common centre
beat the individual raw estimates, even though each raw estimate is unbiased. Their
famous worked example is batting averages, which makes it the right citation for a
sports paper. \S\ref{sec:shrink} applies the idea to per-player draft-position spread,
with the prior fitted from the same pool --- ``empirical Bayes''.

\bibitem{balakrishnan1992}
N.~Balakrishnan (ed.).
\emph{Handbook of the Logistic Distribution.}
Marcel Dekker, 1992.
Standard reference for the logistic distribution: the moment generating function
$\Gamma(1-t)\Gamma(1+t) = \pi t/\sin\pi t$ used in Proposition~\ref{prop:logisticvar},
the variance $\pi^2 s^2/3$, and the excess kurtosis of $6/5$ quoted in
\S\ref{sec:survival}. Any standard text on continuous univariate distributions carries
the same results; the MGF identity itself is Euler's reflection formula.

\bibitem{gneiting2007}
T.~Gneiting and A.~E. Raftery.
\emph{Strictly proper scoring rules, prediction, and estimation.}
Journal of the American Statistical Association, 102(477):359--378, 2007.
The modern reference on propriety. Theorem~\ref{thm:proper} proves the two special cases
this project uses directly rather than invoking the general theory, but this is where
the general theory lives, including why the Brier and logarithmic scores are essentially
the only two anyone should reach for first.

\bibitem{cox1972}
D.~R. Cox.
\emph{Regression models and life-tables.}
Journal of the Royal Statistical Society, Series B, 34(2):187--220, 1972.
The origin of the proportional-hazards idea used in \S\ref{sec:tilt}: multiplying every
subject's hazard by a common factor. Cox's setting is covariate regression on survival
times and the machinery there is far heavier than anything here; what this project
borrows is the single structural observation that scaling cumulative hazard by $c$ is
raising survival to the power $c$, and that this preserves ordering and both fixed
points.

\end{thebibliography}

\end{document}
